\documentclass[aps,pra,twocolumn,showpacs,superscriptaddress,groupedaddress]{revtex4-2}
\usepackage[utf8]{inputenc}
\usepackage{amsmath, amsthm, amssymb, amsfonts,enumitem,listings}
\usepackage[makeroom]{cancel}
\usepackage{graphicx}
\usepackage{bm}
\usepackage{dsfont}
\usepackage{mathtools}
\usepackage{booktabs}
\usepackage{thm-restate}
\usepackage{comment}
\usepackage{pgfplotstable}
\pgfplotsset{compat=1.18}
\usepackage{pgfplots}
\usepgfplotslibrary{groupplots,statistics}
\usepackage{ytableau}
\usepackage[dvipsnames]{xcolor}
\usepackage[colorlinks=true, linkcolor=blue, urlcolor=black,citecolor=blue]{hyperref}
\usepackage{orcidlink}
\DeclareMathOperator{\diag}{diag}
\DeclareMathOperator{\sign}{sign}
\DeclareMathOperator{\imm}{Imm}
\usepackage{physics}
\DeclareMathOperator{\per}{per}
\DeclareMathOperator{\Wg}{Wg}
\DeclareMathOperator{\wg}{wg}
\DeclareMathOperator{\spec}{spec}
\DeclareMathOperator{\Haar}{Haar}
% \DeclareMathOperator{\tr}{tr}
\DeclareMathOperator{\SEP}{SEP}
\DeclareMathOperator{\sep}{sep}
\DeclareMathOperator{\ent}{ent}
% \DeclareMathOperator{\rank}{rank}
\DeclareMathOperator{\id}{id}
\DeclareMathOperator{\Var}{Var}
\DeclareMathOperator{\height}{h}
\newcommand{\GHZ}{{\rm{GHZ}}}
% \newcommand{\bra}[1]{\mathinner{\langle #1|}}
% \newcommand{\ket}[1]{\mathinner{|#1\rangle}}
% \newcommand{\braket}[2]{\mathinner{\langle #1|#2\rangle}}
% \newcommand{\dyad}[1]{| #1\rangle \langle #1|}
\newcommand{\ot}[0]{\otimes}
\newcommand{\one}[0]{\mathds{1}}
\renewcommand{\a}{\alpha}
\renewcommand{\b}{\beta}
\newcommand{\vv}{\ket{v_1} \ot \dots \ot \ket{v_n}}
\newcommand{\vvv}{\dyad{v_1} \ot \cdots \ot \dyad{v_n}}
\newcommand{\pdim}{2^{\lfloor n/2\rfloor}}
\newcommand{\cdn}{(\C^d)^{\otimes n}}
\newcommand{\cdm}{(\C^d)^{\otimes m}}
\newcommand{\instr}[1]{{\color{gray} instr: #1}}
\newcommand{\albert}[1]{{\color{blue} Albert: #1}}
\newcommand{\idk}[1]{{\color{red} #1}}
\definecolor{antongreen}{HTML}{229922}
\newcommand{\anton}[1]{{\color{antongreen} Anton: #1}}
\newcommand{\jofre}[1]{{\color{orange} Jofre: #1}}
\newcommand{\naga}[1]{{\color{red} Naga: #1}}
\newcommand{\gerard}[1]{{\color{brown} Gerard: #1}}
\definecolor{claudeorange}{HTML}{D97757}
\newcommand{\claude}[1]{{\color{claudeorange} Claude: beep boop, #1}}
\newcommand{\fix}[1]{{\color{gray}\footnotesize\,[was: #1]}}
\newcommand{\N}{\mathds{N}}
\newcommand{\R}{\mathds{R}}
\newcommand{\C}{\mathds{C}}
\newcommand{\WW}{\mathcal{W}}
\newcommand*{\bfrac}[2]{\genfrac{\lbrace}{\rbrace}{0pt}{}{#1}{#2}}
\newcommand{\E}{\mathcal{E}}
\DeclarePairedDelimiter\floor{\lfloor}{\rfloor}
\DeclarePairedDelimiter\ceil{\lceil}{\rceil}
\newcommand{\cdk}{(\C^d)^{\otimes k}}
\newcommand{\cd}{\C^d}
\newcommand{\cdx}[1]{(\C^d)^{\otimes #1}}
\newcommand{\I}{\one}

\newcommand{\og}[1]{{\color{orange} #1}}

\DeclareUnicodeCharacter{202C}{\^{i}}
\newcommand{\Om}[1]{\small $\omega_{#1}$}
\newcommand{\De}[1]{$\Delta_{#1}$}
\newcommand{\Ga}[1]{$\Lambda_{4,#1}$}
\makeatletter
\newcommand{\vast}{\bBigg@{4}}
\newcommand{\Vast}{\bBigg@{5}}
\makeatother
\newtheorem{definition}    {Definition}
\newtheorem{theorem}    {Theorem}
\newtheorem{cor}    {Corollary}
\newtheorem{proposition}[theorem]{Proposition}
\newtheorem{observation}[theorem]{Observation}
\newtheorem{lemma}      [theorem]{Lemma}
\newtheorem{example}    [theorem]{Example}
\newtheorem{result}    [theorem]{Result}
\usepackage{phfcc}
\phfMakeCommentingCommand[initials={NT}]{nt}

\usepackage{soul}

\begin{document}
\title{
Entanglement measures and LOCC conversion criteria from Schur-Weyl duality
% Generic convex roofs hierarchies and LOCC conversions
% Entanglement of formation and LOCC conversion from Schur-Weyl duality
}

\author{Anton van Cleemput$^{\orcidlink{0009-0005-3744-3774}}$}
\author{Otfried Gühne$^{\orcidlink{0000-0002-6033-0867}}$}
\author{Albert Rico$^{\orcidlink{0000-0001-8211-499X}}$}

\affiliation{Naturwissenschaftlich-Technische Fakult\"{a}t, Universit\"{a}t Siegen, Walter-Flex-Stra\ss e 3, 57068 Siegen, Germany}



\date{\today}
\begin{abstract}
% The goal of this paper is to impress Marco Túlio Quintino so Anton gets a good grade for his defense. (And a box of chocolate?)
The characterization of possible transformations between quantum states and the quantification of entanglement are central for understanding quantum resources. Explicit results are, however, scarce since they require complicated optimizations over transformation protocols or state decompositions. We introduce a method to optimize functions of the
entanglement spectrum over mixed-state decompositions with arbitrary accuracy. Our results are based on a convergent semidefinite programming hierarchy of optimizations over entanglement spectra using a natural polynomial basis and Schur-Weyl duality. Our approach can also be used 
to determine non-polynomial measures, such as the entanglement of formation,
and to characterize the dimensionality of entanglement. 
Moreover, extracting the decomposition allows us to formulate  
criteria to determine state transformations under local operations
and classical communication (LOCC). In particular, 
we provide a hierarchy of certificates of LOCC non-convertibility 
from pure to mixed states in any dimension and a heuristic approach 
for LOCC convertibility.
\end{abstract}

\maketitle
\section{Introduction}
Composite quantum systems are entangled when they can not be described solely by their parts~\cite{epr}. Operationally, entanglement is a resource for quantum communication~\cite{distributed, eckert} and computation~\cite{one-way, shor} that cannot increase under local operations and classical communications (LOCC) between distant parties.
In isolated two-particle quantum systems, described by pure states, entanglement and LOCC transformations are completely characterized by 
the Schmidt coefficients or the eigenvalues of the marginals~\cite{Nielsen1999MajLOCC, Vidal_2000}, also called 
the entanglement spectrum. Therefore, learning this spectrum is 
necessary and sufficient to understand the entanglement of the 
system and the possible entanglement transformations between 
states.

In practice, however, pure states become mixed due to decoherence with the environment~\cite{Zurek2003,BreuerPetruccione2002,NielsenChuang2010}. When this occurs, the spectrum of the marginals alone is not informative of the entanglement between them, because a mixed state admits multiple descriptions as a probabilistic ensemble of pure states. Consequently, 
one needs to account for all such decompositions. This leads to an intractable optimization required to evaluate entanglement measures~\cite{Bennett1996EoF, Uhlmann1998, Uhlmann2010Roofs, Horodecki2009RMP}, even when full tomography is available.
%~\cite{gurvits2003classicaldeterministiccomplexityedmonds,Huang2014NPhard}.
Existing methods address each specific entanglement measure independently~\cite{Wootters1998, TerhalVollbrecht2000, VollbrechtWerner2001, Mintert2004Concurrence, Chen2005Concurrence, Chen2005EoF, Guhne2007Estimating, Eisert2007QuantWitness}. Moreover, while lower bounds can be found for polynomial measures~\cite{Toth2015CroofTP}, much less is known for the more general case~\cite{PlenioVirmani2007}. As a result, it is still a remarkable challenge to determine when entanglement transformations are possible with LOCC~\cite{Gour2005InfiniteConditions, Chitambar2014LOCC, PhysRevX.8.031020}. 

Here we aim to provide a systematic method to learn the entanglement properties of mixed bipartite quantum states. 
To this end we use the Schur-Weyl decomposition of the Hilbert space of multiple copies, which allows one
to learn entanglement spectra of pure states \cite{Alicki1988, Keyl2001EstimatingSpectrum} to estimate any polynomial entanglement measure for mixed states with a complete polynomial basis. This basis is particularly suited to detect entanglement dimensionality. We find upper and lower converging bounds, which determine, for instance, the amount of entanglement of the main families of bound entangled states. Using semidefinite and linear programming, we extend them to a convergent hierarchy to find any possible entanglement measure of mixed states. 
The combination of these bounds leads us to a hierarchy of certificates to determine LOCC convertibility from pure to mixed and a heuristic approach towards convertibility.

\iffalse

    We build on the work of Ref.~\cite{Toth2015CroofTP}, which relaxes the set of states $\omega_k = \sum_i p_i \dyad{\psi_i}^{\otimes k}$ to obtain SDP lower bounds on the convex roof of the linear entropy, the Schmidt number and the three-tangle. We extend their work in several directions:
    \begin{itemize}
        \item We bound the convex roof of arbitrary symmetric polynomials of the entanglement spectrum via the Schur basis and isotypic projectors.
        \item We prove convergence of the hierarchy using the quantum de Finetti theorem.
        \item We reduce the cost of the SDP by reducing it to the symmetric subspace of the support of $\rho$ and optimizing over real variables when the state is real.
        \item We introduce a cheaper hierarchy, which we use to compute bounds in larger systems.
        \item We embed an LP to bound non-polynomial measures such as the entanglement of formation and Vidal's monotones, and to check LOCC convertibility.
        \item We show how to extract decompositions from the result of the SDP, which give upper bounds in the general case and certificates of convertibility in the LOCC problem, complementing the certificates of non-convertibility given by the SDP.
    \end{itemize}
\fi

\section{Entanglement spectrum}

Consider a bipartite pure state $\ket{\psi}_{AB} \in \C^{d_A} \otimes \C^{d_B} \cong \C^D$ with $D = d_Ad_B$. Such a state is separable if it can be written as $\ket{\psi}=\ket{a}\otimes\ket{b}$ and entangled otherwise. More generally, any pure bipartite state admits a \emph{Schmidt decomposition}
\begin{equation}\label{eq:SchmidtDec}
    \ket{\psi}_{AB} = \sum_{i=1}^r \sqrt{\a_i}\ket{a_i}\otimes\ket{b_i},
\end{equation}
where $\a_1\geq\a_2\geq\dots\geq\a_r>0$ and the sets $\{\ket{a_i}\}_i$ and $\{\ket{b_i}\}_i$ are orthonormal. The spectrum $\a=(\a_1,\dots,\a_{d_A})$ of the reduced density matrix $\rho_A=\tr_B\dyad{\psi}$ (with $\a_i=0$ for $i>r$) is called the {\em entanglement spectrum} of $\ket{\psi}_{AB}$, and the number $r$ of nonzero coefficients is known as the {\em Schmidt rank}, denoted $\mathrm{SR}(\psi)$.
By Nielsen's theorem~\cite{Nielsen1999MajLOCC}, a bipartite pure state $\ket{\psi}$ with entanglement spectrum $\a(\psi)$ can be converted by LOCC into another bipartite pure state $\ket{\phi}$ with entanglement spectrum $\a(\phi)$ if and only if $E_j(\psi)\geq E_j(\phi)$ for all $j=1,\dots,d_A-1$, where
\begin{equation}\label{eq:Vidal}
    E_j(\psi)=1-\sum_{i=1}^{j} \a_i(\psi),
\end{equation}
as defined in Ref.~\cite{Vidal_1999}. In this sense, the quantities $E_j(\psi)$ form a complete set of entanglement monotones, and all information about the entanglement of each state is encoded in the entanglement spectrum. Every pure-state entanglement monotone is a concave, unitarily invariant function of $\rho_A$~\cite{Vidal_2000}.

A pure-state entanglement measure $E$ can be extended to mixed states via the convex roof extension,
\begin{equation}\label{eq:CRoof}
    \widehat{E}(\rho) = \min_{\mathcal{D}(\rho)} \sum_i p_i E(\psi_i),
\end{equation}
where the minimization is taken over all decompositions
$\mathcal{D}(\rho)=\{p_i,\ket{\psi_i}\}$ of $\rho$ such that $\rho = \sum_i p_i \dyad{\psi_i}$. However, this optimization is generally intractable, even if the state is known.
%full tomography is %available~\cite{gurvits2003classica%deterministiccomplexityedmonds}.

\iffalse
    \begin{table}[]
        \centering
        \begin{tabular}{c|ccc}
             & Pure states & Mixed states  \\
             \hline
            with tomography & SVD & This work \\
            without tomography & WSS & Also this work \\
        \end{tabular}
        \caption{Methods to access the average marginal spectrum of a quantum state in different settings.}
        \label{tab:placeholder}
    \end{table}
\fi



\section{Decomposition of the
Hilbert space}
Let us now explain facts about the Hilbert space of $k$ copies of a state $\ket{\psi}$, since this is the basis of our approach.
A general $k$-qudit quantum state is a vector in the Hilbert space $\cdk$. We consider two groups acting on this space: the \emph{symmetric group} $S_k$ of permutations $\pi\in S_k$ of $k$ elements, acting as
\begin{align}
    \pi (\ket{\psi_1} \ot \dots \ot \ket{\psi_k}) &= \ket{\psi_{\pi^{-1}(1)}} \ot \dots \ot \ket{\psi_{\pi^{-1}(k)}},
\end{align}
and the \emph{general linear group} $GL_d$ of invertible $d\times d$ matrices $g\in GL_d$, acting as
\begin{align}
g(\ket{\psi_1} \ot \dots \ot \ket{\psi_k}) &=g\ket{\psi_1} \ot \dots \ot g\ket{\psi_k}.
\end{align}
These two actions commute,
\begin{equation}\label{eq:PerUniCommute}
    \pi \, g = g \, \pi.
\end{equation}
The so-called \emph{Schur--Weyl duality}~\cite{fulton2013representation} states that under the action of these two groups, the Hilbert space $\cdk$ decomposes as
\begin{equation}
    \cdk = \bigoplus_{\substack{\lambda\vdash k \\ \ell(\lambda)\leq d}} \mathcal{S}^\lambda_k \ot \mathcal{U}^\lambda_d
\end{equation}
into irreducible subspaces $\mathcal{S}^\lambda_k$ and $\mathcal{U}^\lambda_d$, which carry the irreducible representations
of $S_k$ and $GL_d$, respectively, indexed by partitions $\lambda$. The dimensions of $\mathcal{S}^\lambda_k$ and $\mathcal{U}^\lambda_d$ are denoted as $f^\lambda$ and $s_\lambda(1,\dots,1)$, where $s_\lambda$ is the so-called Schur polynomial (see also below). Each partition $\lambda$ is represented as a Young diagram: a non-increasing vector of positive integers summing to $k$, whose length we denote by $\ell(\lambda)$.  We write $\lambda \vdash k$ to denote all the partitions $\lambda$ of $k$.

To each Young diagram $\lambda$ we can associate a so-called {\em isotypic projector} $\Pi^\lambda$ onto the subspace $\mathcal{S}^\lambda_k \ot \mathcal{U}^\lambda_d$. This projector is given by
\begin{align}
    \Pi^\lambda = \frac{f^\lambda}{k!}\sum_{\pi \in S_k} \chi^\lambda(\pi)\, \pi,
\end{align}
where $\chi^\lambda(\pi)$ is the character of the cycle type of $\pi$ in the irreducible representation $\lambda$. It follows that the isotypic projectors commute with both permutations and tensor powers of matrices,
\begin{equation}
    \pi\,\Pi^\lambda = \Pi^\lambda\,\pi\,,\quad g\,\Pi^\lambda = \Pi^\lambda\,g\,.
\end{equation}

\section{Polynomial optimization over mixed-state decompositions}
The spectrum of a reduced state $\rho_A$ can be estimated via weak Schur sampling on $k$ copies of $\rho_A$~\cite{Alicki1988,Keyl2001EstimatingSpectrum}: the probability of measuring a Young diagram $\lambda \vdash k$ is given by
\begin{equation}
    \tr(\Pi^\lambda \rho_A^{\otimes k}) = f^\lambda s_\lambda(\spec\rho_A).
\end{equation}
It is known  that the Schur polynomials $s_\lambda$ with $\lambda \vdash k$ and $\ell(\lambda) \leq d_A$ form a basis of the homogeneous symmetric polynomials $P$ of degree $k$ of the spectrum of $\rho_A$ \cite{MacdonaldSymmFunc}.
\begin{equation}
    s_\lambda(\operatorname{\spec\rho_A}) = \sum_{\mu \vdash k} K_{\lambda\mu}m_\mu(\spec\rho_A)
\end{equation}
where $m_\mu$ is the monomial symmetric polynomials
\begin{equation}
    m_\mu(\a) = \sum_{\beta} \a_1^{\beta_1}\a_2^{\beta_2}\dots\a_{d_A}^{\beta_{d_A}}
\end{equation}
where the sum runs over all {\em distinct} permutations $\beta$ of $(\mu_1, \mu_2, \dots, \mu_{\ell(\mu)}, 0,\dots, 0) \in \N^{d_A}$, and $K_{\lambda\mu}$ is a constant known as the Kostka number~\cite{MacdonaldSymmFunc}.

Since $\tr(\rho_A) = 1$ implies $\tr(\rho_A)^{l} = 1$, any polynomial of degree $l < k$ can be lifted to a homogeneous polynomial of degree $k$. Thus, any symmetric polynomial of degree at most $k$ of the entanglement spectrum can be expressed as
\begin{equation}
    P(\spec\rho_A) = \sum_{\lambda\vdash k}c_\lambda s_\lambda(\spec\rho_A).
\end{equation}

In order to estimate any entanglement measure of bipartite states given by a symmetric polynomial, we are interested in the convex roof of these polynomials,
\begin{equation}\label{eq:PrhoABcroof}
    \widehat{P}(\rho_{AB}) = \min_{\mathcal{D}(\rho_{AB})} \sum_i p_i \tr(O_P\, \rho_{A,i}^{\otimes k}),
\end{equation}
where $\rho_{A,i}=\tr_B\dyad{\psi_i}$ and
\begin{equation}\label{eq:OP}
    O_P = \sum_{\lambda \vdash k} \frac{c_\lambda}{f^\lambda} \Pi^\lambda
\end{equation}
for any tuple of coefficients $\vec{c}:=(c_{\lambda_1},\dots,c_{\lambda_t})$, where $t$ is the number of partitions of $k$. $\widehat{P}$ is an entanglement monotone if $P$ is concave~\cite{Vidal_2000}. For any $P$, we always have unitary invariance since $P$ only depends on the spectrum.

\section{Complete hierarchy for lower bounds}
Our approach is based on an idea presented in Ref.~\cite{Toth2015CroofTP}, where polynomial entanglement measures have been estimated using 
optimizations over functions of several copies of a state.
Given a bipartite mixed state $\rho_{AB}$ acting on $\C^{d_A} \otimes \C^{d_B}$, to any decomposition in $\mathcal{D}(\rho_{AB})$ we associate the state
\begin{equation}\label{eq:omega}
    \omega_k = \sum_i p_i \dyad{\psi_i}^{\otimes k}
\end{equation}
acting on $(\C^{d_A} \otimes \C^{d_B})^{\otimes k}$.
%Then, for an observable $O_m$ acting on $m$ copies of a pure state, its convex roof can be written as
%\begin{equation}
%    \min_{\mathcal{D}(\rho)} \tr(O_m \otimes \I_{k-m} \,\omega_k)
%\end{equation}
%for any $k \geq m$.
%Since the set of states $\omega_k$ is intractable. We compute this convex roof via the following SDP relaxation.
The core idea is that a polynomial entanglement measure of degree not larger than $k$ can be written as a linear function on $\omega_k$. Then, optimizing over all decompositions amounts to optimizing over all $\omega_k$ with
the constraint that $\omega_k$ is symmetric and fully separable with respect to the $k$ copies. The separability constraint is difficult to implement, but it can be treated by a hierarchy of semidefinite approximations, similar to the separability criterion of symmetric extensions
\cite{doherty2002}.


Explicitly, we can compute a sequence
of lower bounds on Eq.~\eqref{eq:PrhoABcroof} as $\widehat{P}(\rho)\geq P_{{\rm LB}_k}(\rho)$ with the following semidefinite programming (SDP) relaxation,
\begin{align}
    P_{{\rm LB}_k}(\rho)=\min_{\Omega_k} \ &\tr\!\left(O_P^A \otimes \I^B \, \Omega_k\right)\label{eq:SDPcroof1}\\
    \text{s.t. } & \Omega_k \succeq 0, \nonumber\\
    & \Pi_\text{sym}^{AB,k} \, \Omega_k = \Omega_k\,\Pi_\text{sym}^{AB,k}  = \Omega_k,\nonumber\\
    & \tr_{2,\dots,k} (\Omega_k) = \rho, \nonumber\\
    & \Omega_k^{T_l} \succeq 0 \quad  \forall \, l \in \{1, \dots, \lfloor \tfrac{k}{2}\rfloor\}.\nonumber
\end{align}
Here, $O_P^A$ acts on $A_1,\dots,A_k$. When $\deg(P) = m < k$, $P$ can be lifted to degree $k$, or, equivalently evaluated on $A_1,\dots,A_m$ as both give the same bound (Lemma~\ref{lem:lifting}). $\Pi_\text{sym}^{AB,k}$ is the projector onto the $k$-copy symmetric subspace of $(\C^{d_A}\otimes\C^{d_B})^{\otimes k}$ across copies, and $T_l$ denotes the partial transpose on the first $l$ copies. Since every $\omega_k$ is in the feasible set of this SDP at every level $k$, it gives a lower bound on the convex roof. This yields an SDP hierarchy of $k$-th order lower bounds. This hierarchy is complete, as shown in Appendix~\ref{app:ProofCongNonExtHier}.
\begin{result}\label{thm:ConvergenceSDPcroof}
The bound of Eq.~\eqref{eq:SDPcroof1} is saturated asymptotically, $\lim_{k\rightarrow\infty} P_{{\rm LB}_k}(\rho) = \widehat{P}(\rho)$.
\end{result}
The proof sketch is as follows: by the de Finetti theorem for Bose-symmetric states~\cite{caves2002unknown,Christandl_2007}, the $m$-copy marginal of any feasible $\Omega_k$ is asymptotically close to a mixture $\sum p_i \dyad{\phi_i}^{\otimes m}$ whose first system marginal is asymptotically close to $\rho$ as $k$ increases. By continuity of the convex roof, $P_{{\rm LB}_k} \geq \widehat{P}(\rho) - \delta_k$ where $\lim_{k\to \infty}\delta_k=0$. Moreover, the bound is nondecreasing in $k$ (Appendix~\ref{app:monotonicity}).

Since $\Omega_k$ is supported on the symmetric subspace $\text{Sym}^k(\C^{d_A}\otimes\C^{d_B})\subset(\C^{d_A}\otimes\C^{d_B})^{\otimes k}$, we can exploit this symmetry to reduce the matrix size of $\Omega_k$ from $D^{k}$ to $\dim \text{Sym}^k(\C^{d_A}\otimes\C^{d_B}) = \binom{k+D-1}{k}$. By projecting onto the range of $\rho$, we can further reduce this size to $\binom{k+r-1}{k}$, with $r=\operatorname{rank}\rho$. Moreover, the variable $\Omega_k$ can be taken real for real states. These symmetry reductions are explained in detail in Appendix~\ref{app:SymRed}.

Naturally, using the isotypic polynomial basis allows us to detect the entanglement dimension, i.e. the Schmidt number~\cite{Terhal_2000}, of a target mixed state. To see this, recall that $\Pi^\lambda\ket{e_1}\otimes\dots\otimes\ket{e_k}=0$ if fewer than $\ell(\lambda)$ distinct vectors of an orthonormal basis $\{\ket{e_i}\}$ appear in the product, and therefore $\tr\big(\Pi^\lambda_A\otimes\one_B \dyad{\psi}_{AB}^{\otimes k}\big) = 0$ if and only if $\ell(\lambda) > \operatorname{SR}(\psi_{AB})$~\cite{Christandl_2005,lovitz2026TreeTNtest}.
Therefore, if $P_{{\rm LB}_k}\neq 0$ up to the solver precision, then no pure-state decomposition of $\rho$ exists in which all states have Schmidt rank lower than $\ell(\lambda)$. This is formally described as follows:
\begin{result}
   If $P_{{\rm LB}_k}> 0$ is certified for $O_P=\Pi^\lambda$ with $\lambda\vdash k$, then $\rho$ has Schmidt number at least $\ell(\lambda)$.
\end{result}
As a simple example, we test this Schmidt number hierarchy on the state $ \rho = p \dyad{\phi^+} + (1-p)\dyad{01}$ acting on $\C^3 \otimes \C^3$, with $\ket{\phi^+}=\frac{1}{\sqrt3}\sum_{i=0}^{2}\ket{ii}$, studied in Ref.~\cite{Tavakoli_2024}, which has Schmidt number $3$ for $p > 0$, and we detect it from $p\approx 0.01$ at level $k=3$ with $\lambda = (1,1,1)$. It is easy to see that the state $\rho$ has Schmidt number three for any $p \geq 0$ , since every pure state in a decomposition lies in the span of $\ket{\phi^+}$ and $\ket{01}$, and any such state  has Schmidt rank three.

\section{Simplified relaxation}

Notice that the isotypic projectors $\Pi^\lambda$ only act on Alice's subsystem copies. This leads us to ask whether the SDP can be simplified. In turn, here we present a simplification of the hierarchy above using a smaller variable that allows us to access larger systems, at the cost of weaker bounds at each level $P_{{\rm LB}_{k,m}} \leq P_{{\rm LB}_m}$ with $P_{{\rm LB}_{k,m}}$ defined in Eq.~\eqref{eq:marginals_SDP} below. For $m \geq k$, define the following extension acting on $(\C^{d_A})^{\otimes m}\otimes(\C^{d_B})^{\otimes k}$ (i.e., $m$ copies of $A$ and $k$ copies of $B$),
\begin{equation}\label{eq:omega_km}
    \omega_{k,m} := \sum_i p_i \dyad{\psi_i}^{\otimes k}\otimes \rho_{A,i}^{\otimes (m-k)},
\end{equation}
which satisfies $\tr_{A_{k+1}, \dots, A_m}(\omega_{k,m})=\omega_k$. We relax this operator with a positive semidefinite variable $\Phi_{k,m}$, and denote $\Phi_{k}:=\tr_{A_{k+1}, \dots, A_m}(\Phi_{k,m})$.
Then we can lower bound $\widehat{P}(\rho)$ as $P_{{\rm LB}_{k,m}}\leq \widehat{P}(\rho)$, where
\begin{align}\label{eq:marginals_SDP}
    P_{{\rm LB}_{k,m}}:=&\min_{\Phi_{k,m}} \tr(O_P^A \otimes \I^B \ \Phi_{k,m})\\
    \text{s.t. } & \Phi_{k,m}\succeq 0, \nonumber\\
    & \Pi_\text{sym}^{AB,k} \, \Phi_{k}\, \Pi_\text{sym}^{AB,k}= \Phi_{k},\nonumber\\
    & \tr_{A_2,\dots,A_m,B_2,\dots,B_k} (\Phi_{k,m}) = \rho_{AB}, \nonumber\\
    & \tr_{B_1,\dots,B_k}(\Phi_{k,m}) \text{ is permutation invariant},\nonumber\\
    & \Phi_{k,m}^{T_S} \succeq 0 \quad  \forall \, S.\nonumber
\end{align}
Here $O_P^A$ is the observable of Eq.~\eqref{eq:OP} for a polynomial of degree $m$, acting on $A_1,\dots,A_m$, and $T_S$ denotes the partial transpose on a set $S$ of registers, where transposing a full copy acts on both $A_i$ and $B_i$. Since $\Phi_{k,m}$ can be assumed invariant under permutations of $A_1B_1, \dots, A_kB_k$ and $A_{k+1}, \dots, A_m$, it suffices to impose partial transpose conditions on the tuples $(i, j)$ where each tuple means take the partial transpose on the first $i$ full-copy registers and the first $j$ marginal registers. One such set of representatives is $\{(i,j):0\leq i\leq k,\ 0\leq j<\tfrac{m-k}{2}\}\ \cup\ \{(i,\tfrac{m-k}{2}):0\leq i\leq\lfloor k/2\rfloor\}$ where the second set is only present for even $m-k$ and the tuple $(0,0)$ is excluded.
For $k\geq\deg P$, the marginal $\Phi_k$ of any feasible $\Phi_{k,m}$ is feasible for Eq.~\eqref{eq:SDPcroof1} at level $k$, so that $P_{{\rm LB}_k}\leq P_{{\rm LB}_{k,m}}$ and this hierarchy is also complete in $k$.

\begin{observation}
   At $k=m$, the simplified relaxation only extends full copies on both Alice and Bob and is equal to full hierarchy from Eq.~\ref{eq:SDPcroof1}. In terms of memory usage, this hierarchy is not necessarily cheaper at level $(k_1, m_1)$ with $k_1< m_1$ than it is at level $(k_2, m_2)$ with $k_2 = m_2$. This depends on the rank of the state it is being run on and the rank of its marginals (See App.~\ref{app:supportReduction}). Typically, level $(k_1, m_1)$ is cheaper for full rank states while $(k_2, m_2)$ is cheaper for low-rank ones.
\end{observation}
We show the memory cost of the hierarchy at different levels for states of different ranks in Fig.~\ref{fig:memory} in the Appendix.

\section{Approximating the optimal decomposition}\label{sec:extractingTheDecomposition}
Our next goal is to extract the decomposition $\{p_i, \ket{\psi_i}\}$ from the optimal solution $\Omega_k$ found in the SDP. This will provide us control of how tight our estimates are. Notice that if the relaxation is tight, we have $\Omega_k = \omega_k$. This is in general not the case, and therefore we write $\Omega_k = \omega_k + \Delta$, where $\Delta$ accounts for the gap of the relaxation and solver noise.
For $k \geq 3$, this decomposition is unique up to phases when the operators $\dyad{\psi_i}$ are linearly independent~\cite{Kruskal1977}. When the relaxation is tight and $\Omega_k = \omega_k$, we can then retrieve it exactly with Jennrich's algorithm~\cite{Harshman1970,Leurgans1993}. Since $\Omega_k$ is in general not of this form, we instead use the PARAFAC (or CANDECOMP) algorithm~\cite{Harshman1970,CarrollChang1970} from the TensorLy package~\cite{tensorly}, initialized with Jennrich's algorithm, to find a decomposition $\{p_j, \ket{\psi_j}\}$ that locally minimizes the Frobenius norm
\begin{equation}
    \Theta_k:=\Big\|\Omega_k - \sum_j p_j \dyad{\psi_j}^{\otimes k}\Big\|_F,
\end{equation}
with $\|X\|_F = \sqrt{\sum_{i,j} |X_{ij}|^2}$. Details are given in Appendix~\ref{app:tensor}.
We obtain a decomposition of the state $\sigma = \sum_j p_j \dyad{\psi_j}$, which is close to, but generally different from $\rho$. We turn it into an exact decomposition of $\rho$ by defining
\begin{equation}\label{eq:repair}
    \ket{\varphi_j} = \rho^{1/2}\sigma^{-1/2}\sqrt{p_j}\ket{\psi_j},
\end{equation}
where $\sigma^{-1/2}$ is taken on the support of $\sigma$. If $\operatorname{supp} \rho \subseteq \operatorname{supp} \sigma$, then $\sum_j \dyad{\varphi_j} = \rho^{1/2}\Pi_\sigma\rho^{1/2} = \rho$, with $\Pi_\sigma$ the projector onto the support of $\sigma$. Hence $\{q_j,\ket{\phi_j}\}$ with $q_j = \|\varphi_j\|^2$ and $\ket{\phi_j}=\ket{\varphi_j}/\|\varphi_j\|$ is a decomposition of $\rho$ and gives a valid upper bound on the convex roof.

As a comparison, we also compute upper bounds by gradient descent over ensemble decompositions, which can be unitarily parametrized -- 
a result first proved by Schrödinger \cite{schrodinger36}, but regularly rediscovered~\cite{jaynes57, gisin89, HJW,  kirkpatrick06} (App.~\ref{app:HJW}). Both upper bounds are comparable for the linear entropy and the Vidal monotone $E_1$, while gradient descent gives tighter upper bounds for the entanglement of formation (Figs.~\ref{fig:EoF}, \ref{fig:namedLinearEntropy} and~\ref{fig:Vidal}). However, the tensor extraction method remains useful in cases where an upper bound cannot be extracted otherwise. The main result concerning optimal decompositions can be summarized as follows.

\begin{result}
    Let $P$ be a symmetric polynomial in pure-state Schmidt coefficients.
    It is possible to extract the decomposition of $\Phi_{k,m}$ optimizing $P$. By doing so, we obtain an upper bound $P_{{\rm UB}_k}(\rho)\geq \widehat{P}(\rho)$ which converges, $\lim_{k\rightarrow\infty}P_{{\rm UB}_k}(\rho)=\widehat{P}(\rho)$.
\end{result}\label{thm:extraction}

The proof is shown in App.~\ref{app:extractionConv}, and for it we use the rank reduction step (App.~\ref{app:tensor}), an instance of Jennrich algorithm~\cite{Harshman1970,Leurgans1993} and the repair step from Eq.~\eqref{eq:repair}. We show the bounds on the convex roof of the linear entropy $S_L(\psi)=1-\tr\rho_A^2$ in Fig.~\ref{fig:linearEntropy}.

% Shared colors and styles of the figures. Every figure reads its data from figures/data/*.csv, written by
% experiments/make_figures.py (comma separated, header row, nan for missing values: pgfplots skips those points).
\definecolor{cA}{HTML}{2a78d6}\definecolor{cB}{HTML}{eb6834}\definecolor{cC}{HTML}{1baf7a}%
\definecolor{cD}{HTML}{4a3aa7}\definecolor{cE}{HTML}{c03b7a}\definecolor{cF}{HTML}{8a6d1f}\definecolor{cGray}{HTML}{8a8984}%
\definecolor{lvA}{HTML}{c6dbf5}\definecolor{lvB}{HTML}{93bbec}\definecolor{lvC}{HTML}{5e98e0}%
\definecolor{lvD}{HTML}{2a78d6}\definecolor{lvE}{HTML}{174a8c}\definecolor{lvF}{HTML}{1baf7a}%
%
% \namedbars{csv file}{ylabel}{title}: lower bounds of the named states at every level, divided by the reference `ref`
% (printed above each state at height `label_y`), with the two upper bounds as markers. Columns: see named_csv in
% make_figures.py. The state positions (`pos`) leave a gap at x = 10 between the PPT and the NPT states.
\providecommand{\namedbar}[3]{% fill, column, offset
  \addplot[ybar, bar width=0.13, bar shift=0pt, draw=none, fill=#1, area legend]
    table[x expr=\thisrow{pos}+#3, y expr=\thisrow{#2}/\thisrow{ref}] {\namedfile};}
\providecommand{\namedbars}[3]{%
  \def\namedfile{#1}%
  \begin{tikzpicture}
    \begin{axis}[width=\textwidth, height=0.29\textwidth, ymin=0, ymax=1.38, xmin=0.3, xmax=14.7,
      ylabel={#2}, ytick={0,0.25,0.5,0.75,1}, xtick=data, xticklabels from table={#1}{state},
      table/col sep=comma, xticklabel style={rotate=35, anchor=north east, font=\scriptsize, inner sep=1pt},
      tick label style={font=\footnotesize}, label style={font=\small}, title={#3}, title style={font=\small,
      yshift=8pt}, clip=false, legend style={font=\scriptsize, at={(0.5,-0.42)}, anchor=north, legend columns=4,
      draw=none, /tikz/every even column/.append style={column sep=6pt}}, legend cell align=left,
      ytick style={draw=none}, xtick style={draw=none}, ymajorgrids, grid style={cGray!25}]
      \addplot[draw=none, forget plot] table[x=pos, y expr=0] {#1};% the ticks (xtick=data)
      \namedbar{lvA}{lb12}{-0.325}\addlegendentry{LB $(1,2)$}
      \namedbar{lvB}{lb13}{-0.195}\addlegendentry{LB $(1,3)$}
      \namedbar{lvC}{lb22}{-0.065}\addlegendentry{LB $(2,2)$}
      \namedbar{lvD}{lb23}{0.065}\addlegendentry{LB $(2,3)$}
      \namedbar{lvE}{lb33}{0.195}\addlegendentry{LB $(3,3)$}
      \namedbar{lvF}{lb24}{0.325}\addlegendentry{LB $(2,4)$}
      \addplot[only marks, mark=diamond*, mark size=2.2pt, cB, mark options={draw=white, line width=0.4pt}]
        table[x expr=\thisrow{pos}-0.14, y expr=\thisrow{ub_extract}/\thisrow{ref}] {#1};
      \addlegendentry{UB: extracted ($k=m=3$)}
      \addplot[only marks, mark=x, mark size=2.4pt, black, thick]
        table[x expr=\thisrow{pos}+0.14, y expr=\thisrow{ub_hjw}/\thisrow{ref}] {#1};
      \addlegendentry{UB: HJW}
      \addplot[only marks, mark=none, nodes near coords={$\pgfplotspointmeta$}, point meta=explicit symbolic,
        nodes near coords style={anchor=south, font=\tiny, inner sep=1pt}, forget plot]
        table[x=pos, y=label_y, meta=ref_label] {#1};
      \draw[cGray] (axis cs:0.3,1) -- (axis cs:14.7,1);
      \draw[cGray, densely dashed] (axis cs:10,0) -- (axis cs:10,1.38);
      \node[font=\scriptsize, anchor=south] at (axis cs:5,1.38) {PPT (bound entangled)};
      \node[font=\scriptsize, anchor=south] at (axis cs:12.5,1.38) {NPT};
    \end{axis}
  \end{tikzpicture}}
%
\begin{figure}
  \centering
  \begin{tikzpicture}
    \begin{axis}[width=\columnwidth, height=0.72\columnwidth, xlabel={$a$}, ylabel={bounds on $\widehat{S_L}$},
      ymin=0, ymax=4.2e-3, scaled y ticks=base 10:3, tick label style={font=\footnotesize}, label style={font=\small},
      legend style={font=\scriptsize, at={(0.5,-0.3)}, anchor=north, draw=none, legend columns=3,
      /tikz/every even column/.append style={column sep=4pt}}, legend cell align=left, table/col sep=comma]
      \def\file{figures/data/linear_entropy_horodecki24.csv}
      \addplot[lvA, thick] table[x=a, y=lb12] {\file};\addlegendentry{LB $(1,2)$}
      \addplot[lvB, thick] table[x=a, y=lb13] {\file};\addlegendentry{LB $(1,3)$}
      \addplot[lvC, thick] table[x=a, y=lb22] {\file};\addlegendentry{LB $(2,2)$}
      \addplot[lvD, thick] table[x=a, y=lb23] {\file};\addlegendentry{LB $(2,3)$}
      \addplot[lvE, thick, densely dashed] table[x=a, y=lb33] {\file};\addlegendentry{LB $(3,3)$}
      \addplot[lvF, very thick, dotted] table[x=a, y=lb24] {\file};\addlegendentry{LB $(2,4)$}
      \addplot[only marks, mark=diamond, mark size=3pt, cB, thick] table[x=a, y=ub_extract] {\file};
      \addlegendentry{UB: extracted, $k=m=3$}
      \addplot[only marks, mark=x, black, thick, mark size=1.8pt] table[x=a, y=ub_hjw] {\file};\addlegendentry{UB: HJW}
    \end{axis}
  \end{tikzpicture}
  \caption{Convex roof of the linear entropy $S_L=1-\tr\rho_A^2$ for the Horodecki $2\times4$ family. Lower bounds are shown at different levels of the hierarchy. The markers show the upper bounds extracted from our method and HJW gradient descent.}
  \label{fig:linearEntropy}
\end{figure}


\section{Non-polynomial measures via Linear Programming}\label{sec:LP}
% Generated by experiments/make_figures.py.
\definecolor{cA}{HTML}{2a78d6}\definecolor{cB}{HTML}{eb6834}\definecolor{cC}{HTML}{1baf7a}%
\definecolor{cD}{HTML}{4a3aa7}\definecolor{cGray}{HTML}{8a8984}%
\definecolor{lvA}{HTML}{c6dbf5}\definecolor{lvB}{HTML}{93bbec}\definecolor{lvC}{HTML}{5e98e0}%
\definecolor{lvD}{HTML}{2a78d6}\definecolor{lvE}{HTML}{174a8c}\definecolor{lvF}{HTML}{1baf7a}%
\begin{figure*}
  \centering
  \begin{tikzpicture}
    \begin{axis}[width=\textwidth, height=0.29\textwidth, ymin=0, ymax=1.38, xmin=0.30000000000000004, xmax=14.7,
      ylabel={$E_{F,\mathrm{LB}}\,/\,$ref}, ytick={0,0.25,0.5,0.75,1}, xtick={1,2,3,4,5,6,7,8,9,11,12,13,14},
      xticklabels={{Tiles UPB},{Pyramid UPB},{GenTiles2 UPB},{Chessboard},{Cross-hatch},{Rank-4 PPT},{Horodecki $3\times3$},{Horodecki $2\times4$},{Piani $4\times4$},{Werner $2\times2$},{Gisin},{MEMS},{Bell + $\ket{01}$}}, xticklabel style={rotate=35, anchor=north east, font=\scriptsize, inner sep=1pt},
      tick label style={font=\footnotesize}, label style={font=\small}, title={Entanglement of formation $E_F$}, title style={font=\small,
      yshift=8pt}, clip=false, legend style={font=\scriptsize, at={(0.5,-0.42)}, anchor=north, legend columns=4, draw=none, /tikz/every even column/.append style={column sep=6pt}}, legend cell align=left,
      ytick style={draw=none}, xtick style={draw=none}, ymajorgrids, grid style={cGray!25}]
            \addplot[ybar, bar width=0.13, bar shift=0pt, draw=none, fill=lvA, area legend] coordinates {(0.675,0.324904) (1.675,0.334656) (2.675,0.303153) (3.675,0.404005) (4.675,0.388144) (5.675,0.431027) (6.675,0.123997) (7.675,0.174286) (8.675,0.381966) (10.675,0.482983) (11.675,0.423665) (12.675,0.562028) (13.675,0.437146)};\addlegendentry{LB $(1,2)$}
      \addplot[ybar, bar width=0.13, bar shift=0pt, draw=none, fill=lvB, area legend] coordinates {(0.805,0.537373) (1.805,0.565058) (2.805,0.430906) (3.805,0.517956) (4.805,0.589725) (5.805,0.559301) (6.805,0.205242) (7.805,0.232661) (8.805,0.496878) (10.805,0.592167) (11.805,0.521884) (12.805,0.661474) (13.805,0.544231)};\addlegendentry{LB $(1,3)$}
      \addplot[ybar, bar width=0.13, bar shift=0pt, draw=none, fill=lvC, area legend] coordinates {(0.935,0.610045) (1.935,0.627155) (2.935,0.580622) (3.935,0.606008) (4.935,0.705062) (5.935,0.686084) (6.935,0.302016) (7.935,0.349271) (8.935,0.809017) (10.935,0.827902) (11.935,0.728126) (12.935,0.886515) (13.935,0.767598)};\addlegendentry{LB $(2,2)$}
      \addplot[ybar, bar width=0.13, bar shift=0pt, draw=none, fill=lvD, area legend] coordinates {(1.065,0.61083) (2.065,0.627155) (3.065,0.588493) (4.065,0.709303) (5.065,0.705062) (6.065,0.687289) (7.065,0.308418) (8.065,0.349271) (9.065,0.890301) (11.065,0.827902) (12.065,0.728126) (13.065,0.886515) (14.065,0.767598)};\addlegendentry{LB $(2,3)$}
      \addplot[ybar, bar width=0.13, bar shift=0pt, draw=none, fill=lvE, area legend] coordinates {(1.195,0.610888) (2.195,0.627155) (3.195,0.589237) (4.195,0.774563) (5.195,0.705062) (6.195,0.68729) (7.195,0.308447) (8.195,0.349271) (9.195,0.943262) (11.195,0.827902) (12.195,0.728126) (13.195,0.886515) (14.195,0.767598)};\addlegendentry{LB $(3,3)$}
      \addplot[ybar, bar width=0.13, bar shift=0pt, draw=none, fill=lvF, area legend] coordinates {(1.325,0.923362) (2.325,0.966867) (3.325,0.85094) (4.325,0.95686) (5.325,0.993487) (6.325,0.809194) (8.325,0.530762) (11.325,0.881576) (12.325,0.940658) (13.325,0.98508) (14.325,0.959485)};\addlegendentry{LB $(2,4)$}
            \addplot[only marks, mark=diamond*, mark size=2.2pt, cB, mark options={draw=white, line width=0.4pt}] coordinates {(0.86,1.01207) (1.86,1) (2.86,1.01258) (3.86,1.11694) (4.86,1) (5.86,1) (6.86,1) (7.86,1.00004) (8.86,1.20691) (10.86,1) (11.86,1.0061) (12.86,1.00195) (13.86,1.00003)};\addlegendentry{UB: extracted ($k=m=3$)}
      \addplot[only marks, mark=x, mark size=2.4pt, black, thick] coordinates {(1.14,1) (2.14,1) (3.14,1) (4.14,1) (5.14,1) (6.14,1.00003) (7.14,1.00708) (8.14,1) (9.14,1) (11.14,1) (12.14,1) (13.14,1) (14.14,1)};\addlegendentry{UB: HJW}
      \draw[cGray] (axis cs:0.30000000000000004,1) -- (axis cs:14.7,1);
      \node[anchor=south, font=\tiny, inner sep=1pt] at (axis cs:1,1.06) {$0.214$};
      \node[anchor=south, font=\tiny, inner sep=1pt] at (axis cs:2,1.06) {$0.233$};
      \node[anchor=south, font=\tiny, inner sep=1pt] at (axis cs:3,1.06) {$0.286$};
      \node[anchor=south, font=\tiny, inner sep=1pt] at (axis cs:4,1.157) {$0.55$};
      \node[anchor=south, font=\tiny, inner sep=1pt] at (axis cs:5,1.06) {$0.355$};
      \node[anchor=south, font=\tiny, inner sep=1pt] at (axis cs:6,1.06) {$6.6\mathrm{e}{\scriptscriptstyle-3}$};
      \node[anchor=south, font=\tiny, inner sep=1pt] at (axis cs:7,1.06) {$6.3\mathrm{e}{\scriptscriptstyle-3}$};
      \node[anchor=south, font=\tiny, inner sep=1pt] at (axis cs:8,1.06) {$0.0112$};
      \node[anchor=south, font=\tiny, inner sep=1pt] at (axis cs:9,1.247) {$1$};
      \node[anchor=south, font=\tiny, inner sep=1pt] at (axis cs:11,1.06) {$0.592^*$};
      \node[anchor=south, font=\tiny, inner sep=1pt] at (axis cs:12,1.06) {$0.395^*$};
      \node[anchor=south, font=\tiny, inner sep=1pt] at (axis cs:13,1.06) {$0.722^*$};
      \node[anchor=south, font=\tiny, inner sep=1pt] at (axis cs:14,1.06) {$0.469^*$};
      
      \draw[cGray, densely dashed] (axis cs:10,0) -- (axis cs:10,1.38);\node[font=\scriptsize, anchor=south] at (axis cs:5,1.38) {PPT (bound entangled)};\node[font=\scriptsize, anchor=south] at (axis cs:12.5,1.38) {NPT};
    \end{axis}
  \end{tikzpicture}
  \caption{Lower bounds on the entanglement of formation for different PPT and NPT states, from the SDP with embedded LP (Sec.~\ref{sec:LP}). Bars: lower bounds from the simplified hierarchy at the levels $(k,m)$ shown in the legend. Markers: upper bounds from the decomposition extracted from the SDP solution at $k=m=3$ (orange) and from HJW gradient descent over decompositions (black). All values are divided by the reference printed above each state: the exact convex roof for two qubits (marked $^*$), and the lowest of the two upper bounds otherwise. The gap closes as the level of the hierarchy increases.
  The level $(2,4)$ is missing for the Horodecki $3\times3$ and the Piani $4\times4$ states, since their rank is too high and the SDP exceeded the memory of our machine.}
  \label{fig:EoF}
\end{figure*}
Non-polynomial measures can also be approximated by polynomials.
In theory, non-polynomial functions can be bounded arbitrarily well by infinite series of polynomials. We define
\begin{equation}\label{eq:FLBexact}
    \widetilde{F}_{{\rm LB}_{k,m}}(\rho)=\sup_{\substack{P\leq F\\ \deg P\leq m}}\ \min_{\Phi_{k,m}}\ \tr\left((O_P^A\otimes\I^B)\,\Phi_{k,m}\right),
\end{equation}
the optimization over symmetric polynomials of degree at most $m$ that lower bound $F$ on all spectra.

In practice, the SDP cannot run at large values of $m$, but even at fixed $m$, we can optimize over symmetric polynomials to obtain strong bounds on the convex roof of non-polynomial measures. Here, we show that finding the best polynomial that lower bounds another function is a problem that can be cast as a linear program (LP) over the coefficients in the basis of Schur polynomials. This allows us to dynamically find the best lower-bounding polynomial for a given state while running the SDP.

To build the LP, we pick multiple spectra $\a=(\a_1, \dots, \a_{d_A})$ at which we compute all the Schur polynomials $s_\lambda(\a)$ of degree $m$ and the input function $F(\a)$. We then require $F(\a) \geq \sum_\lambda c_\lambda s_\lambda(\a)$ at every point. By symmetry of the Schur polynomials, we only need to compute the points where $\a_i \geq \a_j$ for $i < j$.
We get the following optimization problem
\begin{equation}\label{eq:FLBgrid}
    F_{{\rm LB}_{k,m}}(\rho)=\max_{\substack{P(\a_g)\leq F(\a_g)\ \forall g\\ \deg P\leq m}}\ \min_{\Phi_{k,m}}\ \tr\left((O_P^A\otimes\I^B)\,\Phi_{k,m}\right).
\end{equation}
This grid procedure tolerates a small error $\epsilon$, where the polynomial might be larger than the given function between grid points. Since we can retrieve the polynomial as output from the SDP, we can compute this error numerically and subtract it from the output to obtain a lower bound. In practice, to obtain a single minimization problem, we replace the LP by its dual (Appendix~\ref{app:LP}).
\begin{result}\label{thm:non-polyHier} We obtain a complete hierarchy to estimate any entanglement measure on mixed states, in the following sense:
\begin{enumerate}
    \item Let $F$ be a continuous symmetric function of the entanglement spectrum, and let $\widehat{F}(\rho)$ be its convex roof~\eqref{eq:CRoof}. Then $\widetilde F_{{\rm LB}_{k,m}}(\rho)\leq\widehat F(\rho)$ for all level $(k,m)$ with $k\leq m$.
    \item The bound is non-decreasing in $k$ and $m$, and $\lim_{k,m\to\infty}\widetilde F_{{\rm LB}_{k,m}}(\rho)=\widehat F(\rho)$. 
    \item With a finite grid, $F_{{\rm LB}_{k,m}}(\rho)-\epsilon\leq\widehat F(\rho)$, where $\epsilon=\max_\a\left(P^*(\a)-F(\a)\right)$ is the grid error of the polynomial $P^*$ returned by the solver.
\end{enumerate}
\end{result}


The proofs of all statements above are given in Appendix~\ref{app:LPconvergence}. We show the bounds obtained from this SDP for the entanglement of formation $E_F=\widehat{S}$ (Fig.~\ref{fig:EoF}), the convex roof of the von Neumann entropy
\begin{equation}\label{eq:vonNeumann}
    S(\psi) = -\tr(\rho_A \log_2 \rho_A), \quad \rho_A = \tr_B \dyad{\psi}.
\end{equation}

\section{LOCC convertibility}
Finally, let us demonstrate how our approach sheds light on the question of which states can be transformed into others by means of LOCC.
By Ref.~\cite{Jonathan_1999}, a pure state $\ket{\psi}$ can be converted into an ensemble $\{p_i, \ket{\psi_i}\}$ by LOCC if and only if
\begin{equation}\label{eq:JPcriterion}
    E_j(\psi) \geq \sum_i p_i E_j(\psi_i)\qquad \text{for all } j=1,\dots,d_A-1,
\end{equation}
where $E_j$ are the Vidal monotones~\eqref{eq:Vidal}.
Hence, $\ket{\psi}$ can be converted into a mixed state $\rho$ if and only if some decomposition of $\rho$ satisfies all these inequalities simultaneously.

Since all the inequalities must hold for the same decomposition, we combine the LPs of Sec.~\ref{sec:LP} for all the $E_j$ into a single SDP. With a grid of spectra $\{\a_g\}_g$ and the moments $\mu_\lambda(\Phi)=\tr\big[(\tfrac{1}{f^\lambda}\Pi^\lambda\otimes\I^B)\Phi\big]$, we define
\begin{align}\label{eq:SDPlocc}
    s_{k,m}=&\min_{\Phi_{k,m},\,y^{(j)},\,s}\  s\\
    &\text{s.t. } \Phi_{k,m}\ \text{feasible for Eq.~\eqref{eq:marginals_SDP}},\quad y^{(j)}\geq0, \nonumber\\
    &\quad \sum_g y^{(j)}_g\,s_\lambda(\a_g)=\mu_\lambda(\Phi_{k,m})\quad\forall\lambda\vdash m,\ \forall j, \nonumber\\
    &\quad \sum_g y^{(j)}_g\,E_j(\a_g)-E_j(\psi)\leq s\quad\forall j. \nonumber
\end{align}
The optimal value $s_{k,m}$ is the smallest worst-case violation of Eq.~\eqref{eq:JPcriterion} over the relaxed set of decompositions. The dual variables of this SDP provide weights $t_j^*\geq0$ with $\sum_jt^*_j=1$ and polynomials $Q_j^*$ satisfying $Q^*_j\leq t^*_jE_j$ on the grid (Appendix~\ref{app:LOCC}). 
% This allows us to address the problem of LOCC convertibility.
As shown in Sec.~\ref{sec:extractingTheDecomposition}, we can extract the decomposition from the result of the SDP. This leads us to a certificate that certain pure-to-mixed LOCC transformations are possible via Eq.~\eqref{eq:JPcriterion}.
\begin{result}
    If the extracted decomposition $\{p_i, \ket{\psi_i}\}$ satisfies Eq.~\eqref{eq:JPcriterion} for a certain input state $\ket{\psi}$ and target state 
    $\rho$, then they can be transformed with LOCC. 
\end{result}
However, unlike for the bounds for polynomial and non-polynomial measures, it may not be possible to extract the correct decomposition, even for large $k$. This comes from the fact that  the underlying decomposition may have a large number of linearly dependent terms and we may not use Jennrich's algorithm. 

\begin{result}\label{thm:LOCC}
We obtain a complete hierarchy for LOCC non-convertibility in the following sense:
    \begin{enumerate}
        \item 
        Consider $\widetilde{s}_{k,m}$, the exact version of Eq.~\eqref{eq:SDPlocc} with $P_j \leq E_j$ on all spectra.
        If $\widetilde{s}_{k,m} > 0$ at some level $(k,m)$, then $\ket{\psi}$ cannot be converted into $\rho$ by LOCC. Moreover, for any $\ket{\psi}$ and $\rho$, if $\ket{\psi} \overset{\text{LOCC}}{\nrightarrow} \rho$, there is a finite level at which $\widetilde{s}_{k,m} > 0$.
    \item With a finite grid, let $\epsilon_j=\max_\a\left(Q^*_j(\a)-t^*_jE_j(\a)\right)$ be the grid errors of the polynomials returned by the solver.
    If $s_{k,m}>\sum_j\epsilon_j$ for some $(k,m)$, then $\ket{\psi}$ cannot be converted into $\rho$ by LOCC.
    \end{enumerate}
\end{result}
The proof is given in Appendix~\ref{app:LOCC}. For $k\geq 3$, the extraction procedure from Section~\ref{sec:extractingTheDecomposition} also provides a heuristic certificate of convertibility: If the extracted decomposition satisfies Eq.~\eqref{eq:JPcriterion}, then $\ket{\psi}$ can be converted into $\rho$ by LOCC.
% Shared colors and styles of the figures. Every figure reads its data from figures/data/*.csv, written by
% experiments/make_figures.py (comma separated, header row, nan for missing values: pgfplots skips those points).
\definecolor{cA}{HTML}{2a78d6}\definecolor{cB}{HTML}{eb6834}\definecolor{cC}{HTML}{1baf7a}%
\definecolor{cD}{HTML}{4a3aa7}\definecolor{cE}{HTML}{c03b7a}\definecolor{cF}{HTML}{8a6d1f}\definecolor{cGray}{HTML}{8a8984}%
\definecolor{lvA}{HTML}{c6dbf5}\definecolor{lvB}{HTML}{93bbec}\definecolor{lvC}{HTML}{5e98e0}%
\definecolor{lvD}{HTML}{2a78d6}\definecolor{lvE}{HTML}{174a8c}\definecolor{lvF}{HTML}{1baf7a}%
%
% \namedbars{csv file}{ylabel}{title}: lower bounds of the named states at every level, divided by the reference `ref`
% (printed above each state at height `label_y`), with the two upper bounds as markers. Columns: see named_csv in
% make_figures.py. The state positions (`pos`) leave a gap at x = 10 between the PPT and the NPT states.
\providecommand{\namedbar}[3]{% fill, column, offset
  \addplot[ybar, bar width=0.13, bar shift=0pt, draw=none, fill=#1, area legend]
    table[x expr=\thisrow{pos}+#3, y expr=\thisrow{#2}/\thisrow{ref}] {\namedfile};}
\providecommand{\namedbars}[3]{%
  \def\namedfile{#1}%
  \begin{tikzpicture}
    \begin{axis}[width=\textwidth, height=0.29\textwidth, ymin=0, ymax=1.38, xmin=0.3, xmax=14.7,
      ylabel={#2}, ytick={0,0.25,0.5,0.75,1}, xtick=data, xticklabels from table={#1}{state},
      table/col sep=comma, xticklabel style={rotate=35, anchor=north east, font=\scriptsize, inner sep=1pt},
      tick label style={font=\footnotesize}, label style={font=\small}, title={#3}, title style={font=\small,
      yshift=8pt}, clip=false, legend style={font=\scriptsize, at={(0.5,-0.42)}, anchor=north, legend columns=4,
      draw=none, /tikz/every even column/.append style={column sep=6pt}}, legend cell align=left,
      ytick style={draw=none}, xtick style={draw=none}, ymajorgrids, grid style={cGray!25}]
      \addplot[draw=none, forget plot] table[x=pos, y expr=0] {#1};% the ticks (xtick=data)
      \namedbar{lvA}{lb12}{-0.325}\addlegendentry{LB $(1,2)$}
      \namedbar{lvB}{lb13}{-0.195}\addlegendentry{LB $(1,3)$}
      \namedbar{lvC}{lb22}{-0.065}\addlegendentry{LB $(2,2)$}
      \namedbar{lvD}{lb23}{0.065}\addlegendentry{LB $(2,3)$}
      \namedbar{lvE}{lb33}{0.195}\addlegendentry{LB $(3,3)$}
      \namedbar{lvF}{lb24}{0.325}\addlegendentry{LB $(2,4)$}
      \addplot[only marks, mark=diamond*, mark size=2.2pt, cB, mark options={draw=white, line width=0.4pt}]
        table[x expr=\thisrow{pos}-0.14, y expr=\thisrow{ub_extract}/\thisrow{ref}] {#1};
      \addlegendentry{UB: extracted ($k=m=3$)}
      \addplot[only marks, mark=x, mark size=2.4pt, black, thick]
        table[x expr=\thisrow{pos}+0.14, y expr=\thisrow{ub_hjw}/\thisrow{ref}] {#1};
      \addlegendentry{UB: HJW}
      \addplot[only marks, mark=none, nodes near coords={$\pgfplotspointmeta$}, point meta=explicit symbolic,
        nodes near coords style={anchor=south, font=\tiny, inner sep=1pt}, forget plot]
        table[x=pos, y=label_y, meta=ref_label] {#1};
      \draw[cGray] (axis cs:0.3,1) -- (axis cs:14.7,1);
      \draw[cGray, densely dashed] (axis cs:10,0) -- (axis cs:10,1.38);
      \node[font=\scriptsize, anchor=south] at (axis cs:5,1.38) {PPT (bound entangled)};
      \node[font=\scriptsize, anchor=south] at (axis cs:12.5,1.38) {NPT};
    \end{axis}
  \end{tikzpicture}}
%
% Classes (column `class` of locc.csv): 0 psi->rho certified (Theorem LOCC (ii)), 1 psi->rho excluded (Theorem LOCC (i)),
% 2 inconclusive, 3 incomparable (only if SHOW_REVERSE in make_figures.py). 31 x 31 grid, rows ordered by p then theta.
\begin{figure}[tbp]
  \centering
  \begin{tikzpicture}
    \begin{axis}[width=0.92\columnwidth, height=0.8\columnwidth, xlabel={$\theta/\pi$}, ylabel={$p$},
      xmin=-0.00416668, xmax=0.254167,
      ymin=-0.0166665, ymax=1.01667, axis on top,
      colormap={locc}{color(0)=(cA!55) color(1)=(cB!60) color(2)=(cGray!45) color(3)=(cC!60)},
      colormap access=direct,
      tick label style={font=\footnotesize}, label style={font=\small}, name=ax]
      \addplot[matrix plot*, mesh/cols=31, point meta=explicit, forget plot]
        table[col sep=comma, x=theta_over_pi, y=p, meta=class] {figures/data/locc.csv};
    \end{axis}
    \node[anchor=south, font=\scriptsize, inner sep=1pt] at ([yshift=3pt]ax.north) {%
      \begin{tabular}{@{}l@{\hspace{1.2em}}l@{}}
      \tikz\fill[cA!55] (0,0) rectangle (2mm,2mm); $\psi\to\rho$ (decomposition) &
      \tikz\fill[cB!60] (0,0) rectangle (2mm,2mm); $\psi\nrightarrow\rho$ (SDP) \\
      \tikz\fill[cGray!45] (0,0) rectangle (2mm,2mm); inconclusive
      \end{tabular}};
  \end{tikzpicture}
  \caption{LOCC convertibility from the pure state $\ket{\psi}$ with entanglement spectrum
  $(0.6,0.3,0.1)$ to the two-qutrit states
  $\rho(p,\theta)=p\dyad{\Phi_3}+(1-p)\dyad{\chi_\theta}$, with $\ket{\Phi_3}$ maximally entangled and
  $\ket{\chi_\theta}=\cos\theta\ket{00}+\sin\theta\ket{11}$, from the SDP of Eq.~\eqref{eq:LOCCminmaxt} at level $(k,m)=(3,4)$.}
  \label{fig:locc}
\end{figure}
. 


\iffalse
\section{Experiment friendly version}
In the constraints of our SDP we might not want to optimize over a state $\rho$ but might want to optimize over all states that satisfy some constraints. This would be the case in an experiment where full tomography is expensive to compute but we might have access to some measurements and statistics about the underlying state. This is of course possible to implement as a constraint in the SDP and instead of the constraint $\tr_{\neq 1}(\Omega_k) = \rho$, we would ask that $\tr_{\neq 1}(\Omega_k)$ satisfies the given constraints.

\fi

\section{Conclusions}
Using spectral estimation techniques and semidefinite programming (SDP), we have introduced a method to estimate any entanglement measure for mixed states. We do so by introducing a complete hierarchy to approximate any measure over mixed-state decompositions by a polynomial lower bound. The presented hierarchy can be adapted with two parameters $(k,m)$ for a tradeoff between efficiency and accuracy. The bounds are optimized with linear programming over a complete polynomial basis that allows us to detect entanglement dimension in a robust manner. Due to the results of Ref.~\cite{Toth2015CroofTP}, the method can also be used without full tomography, by using expectation value constraints on the SDP.

This method also allows us to address the problem of LOCC transformations from pure to mixed states, with a hierarchy of certificates of non-convertibility and a heuristic approach for convertibility. At level $(3,4)$, this already allows us to determine convertibility for a large set of qutrit-qutrit state pairs (see Fig.~\ref{fig:locc}). This resolves a major barrier at certifying  quantum resources and their convertibility, and demonstrates the applicability of Schur-Weyl duality in this area.  

{\em AI statement}. 
The core ideas, methodology,
proof strategies and the project scope
were designed by the authors. 
%The ideas, connections, project scope and %redirections, methodology and proof strategies %were designed by the authors. 
The numerical results, plots and proof calculations were assisted with AI (Claude Opus 4.8-5.5). The article was entirely written by the authors.

{\em Acknowledgements.} We thank \og{some people} for discussions. 
This research was funded by the Deutsche Forschungsgemeinschaft (DFG, German Research Foundation, project number 563437167), the Sino-German Center for Research Promotion (Project M-0294), and the German Federal Ministry of Research, Technology and Space (Project BeRyQC, Grant No.\ 13N17292).




\newpage
\bigskip



\addcontentsline{toc}{subsection}{Bibliography}
\bibliographystyle{ieeetr}
\bibliography{Bibliography}{}

\begin{widetext}



\appendix



\section{Proof of Result~\ref{thm:ConvergenceSDPcroof}}\label{app:ProofCongNonExtHier}
We consider a polynomial $P$ of degree $m$. At level $k \geq m$ of the hierarchy, we obtain the optimal solution $\Omega_k$, with $P_{{\rm LB}_k}(\rho) = \tr(O_P\, \Omega_m)$ where $\Omega_m = \tr_{m+1,\dots, k}(\Omega_k)$ is the $m$-copy marginal of $\Omega_k$.
By the de Finetti theorem for Bose-symmetric states~\cite[Thm.~II.8]{Christandl_2007}, for every feasible $\Omega_k$, there is an $\omega_m = \int \dyad{\phi}^{\otimes m} d\mu(\phi)$ such that
\begin{equation}
    \left\|\Omega_m - \omega_m\right\|_1 \leq \epsilon
\end{equation}
with $\epsilon = \tfrac{4 D m}{k}$. By Carath\'eodory theorem, since the space is finite-dimensional, we can write $\omega_m = \sum_i p_i \dyad{\phi_i}^{\otimes m}$. Here, $\omega_m$ is associated to a decomposition $\{p_i, \ket{\phi_i}\}_i$ corresponding to a state $\sigma$ and $\tr(O_P \,\omega_m) \geq \widehat{P}(\sigma)$. Since the partial trace is contractive, $\left\|\rho - \sigma\right\|_1 = \left\|\tr_{2,\dots, m}(\Omega_m - \omega_m)\right\|_1 \leq \epsilon$.

Then, by H\"{o}lder inequality \eqref{eq:HolderStep}, by definition of the convex roof \eqref{eq:croofstep} and from the continuity of the convex roof of a continuous function in finite dimensions \eqref{eq:contConv}~\cite[Thm.~6.1]{Bucicovschi2013},
\begin{align}
    \tr(O_P \Omega_m) &\geq \tr(O_P\omega_m) - \epsilon\|O_P\|_\infty \label{eq:HolderStep}\\
    &\geq \widehat{P}(\sigma) - \epsilon\|O_P\|_\infty \label{eq:croofstep}\\
    &\geq \widehat{P}(\rho) - \delta(\epsilon) - \epsilon\|O_P\|_\infty \label{eq:contConv}
\end{align}
with $\delta(\epsilon) \to 0$.

Hence, $\widehat{P}{(\rho)} \geq P_{{\rm LB}_k}(\rho) \geq \widehat{P}{(\rho)} -\delta(\epsilon) - \epsilon\|O_P\|_\infty$, so $P_{{\rm LB}_k}(\rho) \to \widehat{P}(\rho)$ as $k \to \infty$.

\section{Lifting}
\begin{lemma}\label{lem:lifting}
    Let $P$ be a homogeneous symmetric polynomial of degree $m$ with $m < k$. Let $P_k$ be its lift to degree $k$ such that $P_k(\a) = P(\a) \left(\sum_i \a_i \right)^{k-m} = P(\a)$ where $\a$ is the spectrum of a state $\rho_A$.
    Then, for every $\Omega_k$ acting on $(\C^{d_A} \otimes \C^{d_B})^{\otimes k}$, which is permutation invariant across all $k$ copies of $\C^{d_A} \otimes \C^{d_B}$,
    \begin{equation}
        \tr(O_P^{A_1,\dots, A_m} \otimes \I^{A_{m+1}, \dots, A_k} \otimes \I^{B} \Omega_k) = \tr(O_{P_k}^A \otimes \I^B \Omega_k).
    \end{equation}
\end{lemma}
\begin{proof}
    Let $\mathcal{T} = \frac{1}{k!}\sum_{\pi \in S_k} \pi \ O_P \otimes \I^{\otimes (k-m)} \ \pi^\dagger$ be the twirl of $O_P \otimes \I^{\otimes (k-m)}$ over $S_k$. Obviously $\mathcal{T}$ commutes with every permutation in $S_k$, and, since it is a linear combination of permutations, we can write it as a sum of isotypic projectors $\mathcal{T} = \sum_{\lambda \vdash k} a_\lambda \Pi^\lambda$ for some coefficients $a_\lambda$.

    For any $\rho_A$ with spectrum $\a$, since $\rho_A^{\otimes k}$ is permutation invariant,
    \begin{equation}
        \tr(\mathcal{T}\rho_A^{\otimes k}) = \tr(\frac{1}{k!} \sum_{\pi \in S_k} \pi O_P \otimes \I \pi^\dagger \rho_A^{\otimes k})= \tr(O_P \otimes \I \rho_A^{\otimes k}).
    \end{equation}
    Thus $\mathcal{T}$ and $O_{P_k}$ have the same expectation value $P_k(\a)$ on every $\rho_A^{\otimes k}$, and since both are combinations of isotypic projectors and the Schur polynomials are linearly independent, $\mathcal{T} = O_{P_k}$.
    Finally,
    \begin{align}
        &\tr(O_P^{A_1,\dots, A_m} \otimes \I^{A_{m+1}, \dots, A_k} \otimes \I^{B} \Omega_k)\\
        =& \tr( \frac{1}{k!}\sum_{\pi \in S_k} O_P^{A_1,\dots, A_m} \otimes \I^{A_{m+1}, \dots, A_k} \otimes \I^{B} \pi^\dagger \Omega_k \pi)\\
        =& \tr( \frac{1}{k!}\sum_{\pi \in S_k}  \pi O_P^{A_1,\dots, A_m} \otimes \I^{A_{m+1}, \dots, A_k} \otimes \I^{B} \pi^\dagger \Omega_k)\\
        =& \tr( O_{P_k} \otimes \I^{B} \Omega_k)
    \end{align}
\end{proof}

\section{Monotonicity in the level}\label{app:monotonicity}
In the full SDP, at level $k+1$, we obtain the optimal solution $\Omega_{k+1}$. Let $O_k$ be the observable at level $k$ and $O_{k+1}$ be its lift to level $k+1$. Then, $P_{{\rm LB}_{k+1}} = \tr(O_{k+1} \Omega_{k+1})=\tr(O_k \otimes \I \ \Omega_{k+1})$ according to Lemma~\ref{lem:lifting}. At level $k$, $\tr_{A_{k+1}B_{k+1}}(\Omega_{k+1})$ is also in the feasible set -- PPT across copies, permutation invariant, first marginal equal to $\rho$ and PSD -- with $\tr(O_k \tr_{A_{k+1}B_{k+1}}(\Omega_{k+1})) = P_{{\rm LB}_{k+1}}$. Then, $P_{{\rm LB}_{k+1}} \geq P_{{\rm LB}_{k}}$.
We can use the same argument in $k$ and $m$ for the simplified SDP.

\section{Symmetry reductions and optimizations}\label{app:SymRed}
Throughout this appendix, a single copy of $\C^{d_A}\otimes\C^{d_B}$ is identified with $\C^D$, with orthonormal basis $\{\ket{j}\}_{j=0}^{D-1}$.

\subsection{Restriction to the symmetric subspace}
The symmetric subspace $\text{Sym}^k(\C^D)$ has an orthonormal basis indexed by occupation vectors. The occupation vector denoted $\nu(i_1,\dots,i_k)\in\N^D$ of a basis vector $\ket{i_1,\dots,i_k}\in(\C^D)^{\otimes k}$ has entries $\nu_j$ equal to the number of times $j$ appears in $(i_1,\dots,i_k)$. For example, $\nu(0,1,2,0,0,1,2,2)=(3,2,3)$ for $D=3$, $k=8$. We define the isometry $V_k:(\C^D)^{\otimes k}\to\text{Sym}^k(\C^D)$ as
\begin{equation}
    V_k = \sum_{i_1, \dots, i_k=0}^{D-1} \binom{k}{\nu(i_1, \dots, i_k)}^{-1/2}\ket{\nu(i_1, \dots, i_k)}\bra{i_1, \dots, i_k},
\end{equation}
where $\binom{k}{\nu}=k!/\prod_j\nu_j!$ is the multinomial coefficient, i.e., the number of basis vectors with occupation vector $\nu$. It has real entries and satisfies $V_kV_k^\dagger=\I_{\text{Sym}^k}$ and $V_k^\dagger V_k=\Pi_{\rm sym}^{k}$. Since every feasible $\Omega_k$ is supported on the symmetric subspace, $\Omega_k=V_k^\dagger\Omega_k^{\rm sym}V_k$ with $\Omega_k^{\rm sym}=V_k\Omega_kV_k^\dagger$, a matrix of size $\binom{k+D-1}{k}$. The objective becomes $\tr\left(V_k(O_P\otimes\I^B)V_k^\dagger\,\Omega_k^{\rm sym}\right)$.

\subsection{First marginal constraint}
Here we wish to encode the condition $\tr_{A_2B_2\dots A_kB_k}(\Omega_k) = \rho$.
We define the occupation vector $\nu_\downarrow$ such that $\ket{\nu_\downarrow} \in \text{Sym}^{k-1}(\C^D)$, and $\ket{\nu_\downarrow + e_j} \in \text{Sym}^k(\C^D)$ has one more element $\ket{j}$, i.e., $(\nu_\downarrow + e_j)_j=(\nu_\downarrow)_j + 1$. Grouping the $\binom{k-1}{\nu_\downarrow}$ tuples $(i_2,\dots,i_k)$ with the same occupation vector and using $\binom{k}{\nu_\downarrow+e_j}=\binom{k-1}{\nu_\downarrow}\frac{k}{\nu_{\downarrow,j}+1}$, we get
\begin{equation}
    \left(\tr_{2,\dots,k}\Omega_k\right)_{jj'}
    =\sum_{\nu_\downarrow}\binom{k-1}{\nu_\downarrow}\frac{\bra{\nu_\downarrow+e_j}\Omega_k^{\rm sym}\ket{\nu_\downarrow+e_{j'}}}{\sqrt{\binom{k}{\nu_\downarrow+e_j}\binom{k}{\nu_\downarrow+e_{j'}}}}
    =\frac1k\tr\!\left(\hat a_{j'}^\dagger\hat a_j\,\Omega_k^{\rm sym}\right),
\end{equation}
with the bosonic creation operators $\hat a_j^\dagger:\text{Sym}^{k-1}(\C^D)\to\text{Sym}^k(\C^D)$,
\begin{equation}
    \hat a_j^\dagger=\sum_{\nu_\downarrow}\sqrt{\nu_{\downarrow,j}+1}\,\ket{\nu_\downarrow+e_j}\bra{\nu_\downarrow}.
\end{equation}
Thus, the marginal constraint becomes the linear constraint $\frac1k\sum_{j,j'}\tr\left(\hat a_{j'}^\dagger\hat a_j\,\Omega_k^{\rm sym}\right)\ket{j}\bra{j'}=\rho$.

\subsection{Partial-transpose constraints}
Thanks to the permutation symmetry, the partial transpose only depends on the size of the set of transposed copies, and since $\Omega_k^{T_{S^c}}=(\Omega_k^{T_S})^T$, a set and its complement give equivalent constraints. We therefore impose the PPT constraints on the first $l$ copies for $l \in \{1, \dots, \lfloor \tfrac{k}{2} \rfloor\}$. We define the isometry
\begin{equation}
    W_l=(V_l\otimes V_{k-l})V_k^\dagger:\ \text{Sym}^k(\C^D)\to\text{Sym}^l(\C^D)\otimes\text{Sym}^{k-l}(\C^D),
\end{equation}
so that $\Omega_k=(V_l\otimes V_{k-l})^\dagger\,W_l\Omega_k^{\rm sym}W_l^\dagger\,(V_l\otimes V_{k-l})$. Writing $T_l$ also for the transpose of the factor $\text{Sym}^l(\C^D)$ in the occupation basis, the partial transpose becomes
\begin{align}
    \Omega_k^{T_l} &= \left((V_l\otimes V_{k-l})^\dagger\,W_l\Omega_k^{\rm sym}W_l^\dagger\,(V_l\otimes V_{k-l})\right)^{T_l}\\
    &=\left(V_l^T\otimes V_{k-l}^\dagger\right)\,\left(W_l\Omega_k^{\rm sym}W_l^\dagger\right)^{T_l}\,\left(V_l^*\otimes V_{k-l}\right)\\
    &= (V_l\otimes V_{k-l})^\dagger\,\left(W_l\Omega_k^{\rm sym}W_l^\dagger\right)^{T_l}\,(V_l\otimes V_{k-l}),
\end{align}
since $V_l$ is real, so that $V_l^T=V_l^\dagger$ and $V_l^*=V_l$. Conjugating back with $V_l\otimes V_{k-l}$, the condition becomes
\begin{equation}
    \Omega_k^{T_l}\succeq0\iff\left(W_l\,\Omega_k^{\rm sym}W_l^\dagger\right)^{T_l}\succeq0,\qquad 1\leq l\leq\lfloor k/2\rfloor.
\end{equation}
This is the most expensive constraint of the SDP, since it encodes the state in the larger space $\text{Sym}^l(\C^D)\otimes\text{Sym}^{k-l}(\C^D)$ to support the partial transpose. While these constraints are not necessary for the hierarchy to be complete, they help convergence at finite $k$, as shown in Ref.~\cite{Navascues2009}.

\subsection{Real variables}
If $\rho$ is real, the SDP can be restricted to real symmetric matrices. If $\Omega_k$ is feasible, so is its entrywise complex conjugate $\Omega_k^*$, since conjugation commutes with partial traces, partial transposes and the real projector $\Pi_{\rm sym}^k$, and $\rho^*=\rho$. Since the isotypic projectors are linear combinations of permutation matrices with real characters, $O_P$ is real, and
\begin{equation}
    \tr\left((O_P^A \otimes \I^B)\, \Omega_k\right) = \tr\left((O_P^A \otimes \I^B)\, \Omega_k^*\right) = \tr\Big((O_P^A \otimes \I^B)\, \tfrac{\Omega_k + \Omega_k^*}{2}\Big).
\end{equation}
By convexity, $\frac12(\Omega_k+\Omega_k^*)$ is feasible, real and has the same objective value, so the optimum is attained by a real matrix. For the true $\omega_k$, this real matrix is $\frac{1}{2}\sum_i p_i \left(\dyad{\psi_i}^{\otimes k} + \dyad{\psi_i^*}^{\otimes k}\right)$, which corresponds to the decomposition $\{\frac{p_i}{2},\ket{\psi_i}\}\cup\{\frac{p_i}{2},\ket{\psi_i^*}\}$ of $\rho$. The same argument applies to $\Phi_{k,m}$ in the simplified relaxation.

\subsection{Projection onto the support of $\rho$}\label{app:supportReduction}
Let $r = \operatorname{rank} \rho$, and let $V_\rho: \C^D \to \C^r$ be the isometry onto the support of $\rho$, given by the eigenvectors of $\rho$ with nonzero eigenvalue, so that $V_\rho V_\rho^\dagger=\I_r$ and $V_\rho^\dagger V_\rho=\Pi_\rho$, the projector onto the support of $\rho$. Every feasible $\Omega_k$ is supported on $(\operatorname{supp}\rho)^{\otimes k}$: since
\begin{equation}
    \tr\left(((\I-\Pi_\rho)\otimes\I)\,\Omega_k\right)=\tr\left((\I-\Pi_\rho)\rho\right)=0
\end{equation}
and $\Omega_k\succeq0$, we have $((\I-\Pi_\rho)\otimes\I)\,\Omega_k=0$, and by permutation symmetry the same holds on every copy. Hence $\Omega_k=\Pi_\rho^{\otimes k}\Omega_k\Pi_\rho^{\otimes k}$, and
\begin{equation}
    \tr\left((O_P^A\otimes\I^B)\,\Omega_k\right)=\tr\Big(\left(V_\rho^{\otimes k}(O_P^A\otimes\I^B)V_\rho^{\dagger\otimes k}\right)\left(V_\rho^{\otimes k}\Omega_kV_\rho^{\dagger\otimes k}\right)\Big).
\end{equation}
We can therefore equivalently optimize over $\Omega_k^\rho=V_\rho^{\otimes k}\Omega_kV_\rho^{\dagger\otimes k}$ on $(\C^r)^{\otimes k}$ with marginal $V_\rho\rho V_\rho^\dagger$. Bose symmetry is preserved, since $V_\rho^{\otimes k}$ commutes with permutations, and the partial-transpose constraints are equivalent, since $(\Omega_k^\rho)^{T_l}$ and $\Omega_k^{T_l}$ are related in both directions by congruences with $V_\rho$ and $V_\rho^*$. Combined with the symmetric subspace, the variable has size $\binom{k+r-1}{k}$. If $\rho$ is real, $V_\rho$ can be chosen real and the real-variable reduction still applies.

The same argument applies to the simplified relaxation: by the permutation invariance of $\tr_B\Phi_{k,m}$, every $A$ register has reduced state $\rho_A=\tr_B(\rho)$, so all $A$ registers can be encoded on the support of $\rho_A$.


\section{Tensor decomposition}
\subsection{Heuristic extraction}\label{app:tensor}
We take the optimal solution $\Phi_{k,m}$ of Eq.~\eqref{eq:marginals_SDP} (or $\Omega_k$ for the full SDP) and trace out the marginal registers, which we cannot use to extract tensor powers, to obtain $\Omega_k$. We work on the support of $\rho$ as in Appendix~\ref{app:supportReduction}: with $r=\operatorname{rank}\rho$ and the isometry $V_\rho:\operatorname{supp}\rho\to\C^r$, we consider $\Omega_\rho=V_\rho^{\otimes k}\Omega_kV_\rho^{\dagger\otimes k}$.

To work with real tensors, let $\{G_a\}_{a=0}^{r^2-1}$ be a Hilbert--Schmidt orthonormal basis of Hermitian operators on $\C^r$, consisting of the normalized identity $G_0=\I/\sqrt r$ and the generalized Gell-Mann matrices. Every pure state can be written as $\dyad{\psi}=\sum_ax_aG_a$ with a real vector $x$ satisfying $x_0=1/\sqrt r$ and $\|x\|_2=1$. We define the real symmetric order-$k$ tensor $\Omega^G_{a_1\dots a_k}=\tr\left((G_{a_1}\otimes\dots\otimes G_{a_k})\,\Omega_\rho\right)$, which for the true $\omega_k$ equals $\sum_ip_i\,x^{(i)\otimes k}$. We run PARAFAC (alternating least squares) on $\Omega^G$ with nonnegative weights, obtaining terms $\{p_i,v_i\}$; we normalize each factor such that $v_{i,0}=1/\sqrt r$, absorbing the scale into $p_i$, and recover the pure states $\ket{\psi_i}$ as the top eigenvectors of $\sum_av_{i,a}G_a$. Finally, we repair the decomposition $\{p_i,\ket{\psi_i}\}$ as described in Sec.~\ref{sec:extractingTheDecomposition}.


The number of terms $n$ fitted by PARAFAC is the rank of $\Omega_\rho$, which equals the number of terms when the $\dyad{\psi_i}^{\otimes k}$ are linearly independent. 
Any decomposition of $\rho$ has at least $r$ terms, otherwise the repair fails. Conversely, every decomposition of $\rho$ consists of states on $\operatorname{supp}\rho\cong\C^r$, whose state space has dimension $r^2$, so by Carath\'eodory theorem the convex roof of a continuous function is attained by a decomposition with at most $r^2$ terms~\cite[Lemma~A-2]{Uhlmann1998}. 

Before the extraction, it helps to reduce the rank of the SDP solution. We solve a second SDP, minimizing $\tr(R_3 \ot \I \,\Phi_{k,m})$ for a random Hermitian $R_3$ acting on the first $3$ copies, over the same feasible set with the additional constraint $\tr\left((O_P\otimes\I^B)\Phi_{k,m}\right)\leq P_{\rm LB}+\delta_k$ for a small $\delta_k\to0$ with $\delta_k \geq \widehat{P}(\rho)-P_{{\rm LB}_k}$. This selects an extreme point of the near-optimal face, which generically has low rank, instead of a mixture of several optimal solutions. 

\subsection{Convergence of the extraction}\label{app:extractionConv}
Given a state $\omega_3 = \sum_i p_i \dyad{\psi_i}^{\otimes 3}$, we can use Jennrich's algorithm to recover uniquely the decomposition $\{p_i, \ket{\psi_i}\}_i$ when the $\dyad{\psi_i}$ are linearly independent. Therefore, we need to select a low-rank extreme point $\Omega_3 = \tr_{A_4B_4\dots A_kB_k}(\Omega_k)$. For this we apply the rank reduction step as shown in App.~\ref{app:tensor}.

Now we need the following observation.
\begin{observation}
    A decomposition $\{q_i, \dyad{\psi_i}\}$ of $\rho$ is extremal in the set of optimal decompositions of $\rho$ with regards to a polynomial $P$, only if $\dyad{\psi_i}$ are linearly independent projectors.
\end{observation}
\begin{proof}
We will do so by contradiction as follows. Given an optimal decomposition $\{q_i, \dyad{\psi_i}\}$ of $\rho$, if the projectors $\dyad{\psi_i}$ are linearly dependent, there exists a nonzero real vector $c$ such that $\sum_i c_i \dyad{\psi_i} = 0$ with $\sum_i c_i=0$.
Assuming that this is the case, for a small $0<t\leq \min_i \tfrac{q_i}{|c_i|}$, let us define two other decompositions $\{q_i + t c_i, \dyad{\psi_i}\}$ and $\{q_i - t c_i, \dyad{\psi_i}\}$ such that their midpoint is the decomposition $\{q_i, \dyad{\psi_i}\}$.
Since $\sum_i (q_i+t c_i) P(\psi_i) \geq \widehat{P}(\rho)=\sum_i q_i P(\psi_i)$ and $\sum_i (q_i- t c_i) P(\psi_i) \geq \widehat{P}(\rho)$, then $\sum_i (q_i+ t c_i) P(\psi_i) = \sum_i (q_i- t c_i) P(\psi_i)= \widehat{P}(\rho)$. 
% \begin{align}
%     \text{Since } 
%     \sum_i q_i P(\psi_i) &=\widehat{P}(\rho), \\
%     \sum_i (q_i+t c_i) P(\psi_i) &\geq \widehat{P}(\rho), \\
%     \sum_i (q_i-t c_i) P(\psi_i) &\geq \widehat{P}(\rho),\\
%     \text{then } \sum_i t c_i P(\psi_i) &\geq 0 \text{ and}\\
%     \sum_i t c_i P(\psi_i) &\leq 0.\\
%     \text{Therefore by combining both bounds we obtain} \sum_i t c_i P(\psi_i) &= 0,\\
%     \text{which means that }\sum_i (q_i+t c_i) P(\psi_i) &= \widehat{P}(\rho), \\
%     \sum_i (q_i-t c_i) P(\psi_i) &= \widehat{P}(\rho),\\
% \end{align}

These decompositions are also optimal and their mid point is the decomposition extracted $\{q_i, \dyad{\psi_i}\}$, which proves that it is not an extreme point. Therefore, any extreme point of the convex set of optimal decompositions must have linearly independent projectors.
\end{proof}

For almost every $R$, the minimizer $\omega_3^*$ of $\tr(R\,\omega_3)$ among the $\omega_3$ of optimal decompositions is unique. The decomposition of $\rho$ giving $\omega_3^*$ has linearly independent projectors. Using the same argument shown above, if it has linearly dependent terms, we can define two decompositions whose associated $\omega_3$ is optimal with midpoint $\omega_3^*$. By uniqueness, both decompositions give $\omega_3^*$. We can simply increase $t$ until a weight vanishes and repeat to end up at a decomposition into linearly independent terms. Therefore, $\omega_3^*$ decomposes into linearly independent terms.

As $k\to\infty$ the rank-reduced $\Omega_3$ converges to $\omega_3(\mu_e)=\sum_ip_i\dyad{\psi_i}^{\otimes3}$, where
$\mu_e$ is the extreme optimal decomposition singled out by $R$, since $\delta_k \geq \widehat{P}(\rho)-P_{{\rm LB}_k}$ leaves $\omega_3(\mu_e)$ feasible during the rank reduction step.
By App.~\ref{app:ProofCongNonExtHier}, every limit point as $k\to \infty$ corresponds to the state $\omega_3$ of an optimal decomposition, and for almost every $R$ the minimizer of $\tr(R\,\cdot)$ on this set is unique and has linearly independent projectors. 
Then, Jennrich's algorithm is continuous at $\omega_3(\mu_e)$, so applied to $\Omega_3$ it gives a decomposition of a state $\sigma$ that is $\epsilon$-close to $\rho$ with $\epsilon\to0$ as $k\to\infty$, which after the repair step has value converging to $\widehat{P}(\rho)$.

This argument does not extend to LOCC since the optimal decompositions do not form a face of the set of decompositions of $\rho$, and the extreme points may not have linearly independent projectors, in which case we cannot use Jennrich's algorithm.

\section{Linear program for non-polynomial functions}\label{app:LP}

\subsection{Construction of the LP and of the single SDP}
We optimize over symmetric polynomials of degree at most $m$. Since polynomials of lower degree can be lifted to degree $m$ (Lemma~\ref{lem:lifting}), we parametrize them as $P(\a)=\sum_{\lambda\vdash m}c_\lambda s_\lambda(\a)$, and the LP variables are the coefficients $\vec c=(c_\lambda)_\lambda$. Given the function $F$, we require $P(\a)\leq F(\a)$ for all spectra $\a$, and enforce this constraint discretely on a grid of spectra $\{\a_g\}_g$, where the values $s_\lambda(\a_g)$ and $F(\a_g)$ are precomputed. By symmetry, it suffices to take grid points with $\a_{g,1}\geq\a_{g,2}\geq\dots\geq\a_{g,d_A}\geq0$ and $\sum_i\a_{g,i}=1$.

Writing $\mu_\lambda(\Phi)=\tr\left((\tfrac{1}{f^\lambda}\Pi^\lambda\otimes\I^B)\Phi\right)$, the objective is $\tr\left((O_P\otimes\I^B)\Phi_{k,m}\right)=\sum_\lambda c_\lambda\,\mu_\lambda(\Phi_{k,m})$, so the bound reads
\begin{equation}
    F_{{\rm LB}_{k,m}}=\max_{\vec c}\ \min_{\Phi_{k,m}}\ \sum_\lambda c_\lambda\,\mu_\lambda(\Phi_{k,m})
    \quad\text{s.t.}\quad \sum_\lambda c_\lambda s_\lambda(\a_g)\leq F(\a_g)\ \ \forall g,
\end{equation}
with $\Phi_{k,m}$ ranging over the feasible set of Eq.~\eqref{eq:marginals_SDP}. Since the objective is bilinear, the set of feasible $\Phi_{k,m}$ is convex and compact, and the set of valid $\vec{c}$ is convex, by Sion minimax theorem~\cite{Sion1958} we can exchange the maximum and the minimum. Then, for fixed $\Phi_{k,m}$, the inner maximization is an LP whose dual is
\begin{equation}
    \min_{y\geq0}\ \sum_g y_g F(\a_g)\quad\text{s.t.}\quad \sum_g y_g\,s_\lambda(\a_g)=\mu_\lambda(\Phi_{k,m})\ \ \forall\lambda .
\end{equation}
The primal is always feasible, since the constant polynomial $P=\min_gF(\a_g)$ is admissible, so strong duality holds. We merge the LP and the SDP to obtain a single SDP,
\begin{align}
    F_{{\rm LB}_{k,m}}=\min_{\Phi_{k,m},\,y}\ & \sum_g y_g F(\a_g)\\
    \text{s.t. }& \Phi_{k,m}\ \text{feasible for Eq.~\eqref{eq:marginals_SDP}},\quad y\geq0, \nonumber\\
    & \sum_g y_g\,s_\lambda(\a_g)=\tr\left((\tfrac{1}{f^\lambda}\Pi^\lambda\otimes\I^B)\,\Phi_{k,m}\right)\quad\forall\lambda\vdash m. \nonumber
\end{align}
The coefficients $c_\lambda$ of the optimal polynomial are recovered as the dual variables of the equality constraints. Since the grid is discrete, the optimal polynomial $P^*$ may exceed $F$ between grid points. We therefore report $F_{{\rm LB}_{k,m}}-\epsilon$, where $\epsilon=\max_\a\left(P^*(\a)-F(\a)\right)$ is estimated from the polynomial returned by the solver. This is a valid lower bound, since $P^*-\epsilon\leq F$ on all spectra. 

\subsection{Convergence of the bound}\label{app:LPconvergence}
Here we prove Result~\ref{thm:non-polyHier} for the bound $\widetilde F_{{\rm LB}_{k,m}}$ of Eq.~\eqref{eq:FLBexact}, in which $P\leq F$ is imposed on all spectra. Monotonicity follows naturally from the facts that admissible polynomials stay admissible at higher level by lifting (Lemma~\ref{lem:lifting}), and that the SDP is monotone (Appendix~\ref{app:monotonicity}).

By the Weierstrass approximation theorem, for any continuous function $F'$ and any $\eta>0$, we can find a polynomial $P'$ such that $|F'(\a)-P'(\a)| < \eta$ for all spectra $\a$. Then, by averaging over permutations of the variables, for any continuous symmetric function $F$, we can find a symmetric polynomial $P$ such that $|F(\a)-P(\a)| < \eta$ for all spectra $\a$.
Shifting $P$ by $\eta$, for any $F$ and any $\eta>0$, there exists a symmetric polynomial $P$ such that $P(\a) < F(\a)$ and $F(\a)-P(\a) < 2\eta$ for all spectra $\a$. This $P$ is admissible for $m\geq k\geq\deg P$, so by Result~\ref{thm:ConvergenceSDPcroof}, $\lim_{k,m\to\infty}\widetilde F_{{\rm LB}_{k,m}}(\rho)\geq\lim_{k,m\to\infty}P_{{\rm LB}_k}(\rho)=\widehat P(\rho)\geq\widehat F(\rho)-2\eta$. Conversely, $P\leq F$ implies $\widehat P(\rho)\leq\widehat F(\rho)$ for every admissible $P$, so $\widetilde F_{{\rm LB}_{k,m}}(\rho)\leq\widehat F(\rho)$. Since $\eta$ is arbitrary,
\begin{equation}
    \lim_{k,m\to\infty}\widetilde F_{{\rm LB}_{k,m}}(\rho)=\widehat F(\rho).
\end{equation}

With a finite grid, the set of admissible polynomials is larger, so $F_{{\rm LB}_{k,m}}\geq\widetilde F_{{\rm LB}_{k,m}}$, and $F_{{\rm LB}_{k,m}}-\epsilon\leq\widehat F(\rho)$, where $\epsilon$ is the grid error of the returned polynomial. The uncorrected value $F_{{\rm LB}_{k,m}}$ is nondecreasing in the level by the same argument as above, but the corrected value $F_{{\rm LB}_{k,m}}-\epsilon$ need not be, since the grid error can change between levels. Since every corrected value is a valid lower bound, their running maximum over the levels is a valid and nondecreasing lower bound. In our numerics, increasing the level yields tighter bounds.

\section{LOCC feasibility SDP}\label{app:LOCC}
\subsection{Construction of the SDP}
Let $\ket{\psi}$ be the initial pure state and $\rho$ the target mixed state. By Ref.~\cite{Jonathan_1999}, $\ket{\psi}$ can be converted into $\rho$ by LOCC if and only if there is a decomposition $\{p_i,\ket{\psi_i}\}$ of $\rho$ such that
\begin{equation}\label{eq:LOCCcriterion}
    \sum_i p_i E_j(\psi_i)\leq E_j(\psi)\qquad\text{for all } j=1,\dots,d_A-1 .
\end{equation}
Since all these inequalities must hold for the same decomposition, testing them one monotone at a time is not sufficient. Instead, we combine the LP construction of Appendix~\ref{app:LP} for all the Vidal monotones into a single SDP.

For each $j$, let $\mathcal{P}_j$ be the set of symmetric polynomials $P_j$ of degree at most $m$ that satisfy $P_j(\a_g)\leq E_j(\a_g)$ on all grid points, and let
\begin{equation}
    \epsilon_j(P_j)=\max_{\a}\left(P_j(\a)-E_j(\a)\right)
\end{equation}
be the maximal violation of $P_j$ between grid points. If a decomposition $\{p_i,\ket{\psi_i}\}$ satisfies Eq.~\eqref{eq:LOCCcriterion}, then its operator $\omega_{k,m}$ is feasible for Eq.~\eqref{eq:marginals_SDP} and, for every $P_j\in\mathcal{P}_j$,
\begin{equation}\label{eq:LOCCnecessary}
    \tr\left((O_{P_j}\otimes\I^B)\,\omega_{k,m}\right)=\sum_ip_iP_j(\a_i)\leq\sum_ip_iE_j(\a_i)+\epsilon_j(P_j)\leq E_j(\psi)+\epsilon_j(P_j),
\end{equation}
where $\a_i$ is the entanglement spectrum of $\ket{\psi_i}$. In words, a convertible decomposition can appear to violate the $j$-th inequality by at most the grid error of the polynomial used to test it.

With the moments $\mu_\lambda(\Phi)=\tr\left((\tfrac{1}{f^\lambda}\Pi^\lambda\otimes\I^B)\Phi\right)$ of Appendix~\ref{app:LP}, the SDP of Eq.~\eqref{eq:SDPlocc} reads
\begin{align}\label{eq:LOCCsdp}
    s_{k,m}=\min_{\Phi_{k,m},\,y^{(j)},\,s}\ & s\\
    \text{s.t. }& \Phi_{k,m}\ \text{feasible for Eq.~\eqref{eq:marginals_SDP}},\quad y^{(j)}\geq0\ \ \forall j, \nonumber\\
    & \sum_g y^{(j)}_g\,s_\lambda(\a_g)=\mu_\lambda(\Phi_{k,m})\quad\forall\lambda\vdash m,\ \forall j, \nonumber\\
    & \sum_g y^{(j)}_g\,E_j(\a_g)-E_j(\psi)\leq s\quad\forall j. \nonumber
\end{align}

For fixed $\Phi_{k,m}$, the minimization over each $y^{(j)}$ is the dual LP of Appendix~\ref{app:LP} with $F=E_j$, so by strong duality
\begin{equation}\label{eq:LOCCminmax}
    s_{k,m}=\min_{\Phi_{k,m}}\ \max_{j}\ \max_{P_j\in\mathcal{P}_j}\Big(\tr\left((O_{P_j}\otimes\I^B)\Phi_{k,m}\right)-E_j(\psi)\Big).
\end{equation}
Thus $s_{k,m}$ is the smallest worst-case violation of Eq.~\eqref{eq:LOCCcriterion} over the relaxed set of decompositions, where each violation is measured with the best grid-admissible polynomial. For a single monotone, $s_{k,m}=F_{{\rm LB}_{k,m}}-E_j(\psi)$ with $F=E_j$.

Let $t$ be a probability vector with $t_j \geq 0$ and $\sum_j t_j = 1$. Then,
\begin{align}\label{eq:LOCCminmaxT}
    s_{k,m}=\min_{\Phi_{k,m}} \max_t \max_{P_j\in\mathcal{P}_j}\ \sum_j t_j\Big(\tr\left((O_{P_j}\otimes\I^B)\Phi_{k,m}\right)-E_j(\psi)\Big).
\end{align}
With the substitution $Q_j=t_jP_j$, the constraint $P_j\in\mathcal{P}_j$ becomes $Q_j(\a_g)\leq t_jE_j(\a_g)$ for all $g$, and the objective is linear in $\Phi_{k,m}$ and in $(t,Q_j)$. The feasible set of $\Phi_{k,m}$ is convex and compact, and the set of $(t,Q_j)$ is convex. By Sion minimax theorem~\cite{Sion1958}, the minimum and the maximum can be exchanged:
\begin{align}
    s_{k,m}&=\min_{\Phi_{k,m}}\ \max_{t,\,Q_j}\ \sum_j \Big(\tr\left((O_{Q_j}\otimes\I^B)\Phi_{k,m}\right)-t_jE_j(\psi)\Big)\nonumber\\
    &=\max_{t,\,Q_j}\ \min_{\Phi_{k,m}}\ \sum_j \Big(\tr\left((O_{Q_j}\otimes\I^B)\Phi_{k,m}\right)-t_jE_j(\psi)\Big).\label{eq:LOCCmaxmin}
\end{align}
Let $t^*$ and $Q_j^*$ be the optimal vector $t$ and optimal polynomials, which are given by the dual variables of the SDP~\eqref{eq:LOCCsdp}; for $t^*_j>0$, $Q^*_j=t^*_jP^*_j$ with $P^*_j\in\mathcal{P}_j$.
By Eq.~\eqref{eq:LOCCmaxmin}, they satisfy
\begin{equation}\label{eq:LOCCcertificateIneq}
    \sum_j \Big(\tr\left((O_{Q^*_j}\otimes\I^B)\Phi_{k,m}\right)-t^*_jE_j(\psi)\Big)\geq s_{k,m}
\end{equation}
for every feasible $\Phi_{k,m}$.

Let $\epsilon_j=\max_\a\big(Q^*_j(\a)-t^*_jE_j(\a)\big)$, which equals $t^*_j\,\epsilon_j(P^*_j)$ for $t^*_j>0$. If $\ket{\psi}$ can be converted into $\rho$ by LOCC, some decomposition satisfies Eq.~\eqref{eq:LOCCcriterion}, and its $\omega_{k,m}$ is feasible. As in Eq.~\eqref{eq:LOCCnecessary}, $\tr\big((O_{Q^*_j}\otimes\I^B)\,\omega_{k,m}\big)=\sum_ip_iQ^*_j(\a_i)\leq t^*_jE_j(\psi)+\epsilon_j$, so Eq.~\eqref{eq:LOCCcertificateIneq} gives
\begin{equation}
    s_{k,m} \leq \sum_j \Big(\tr\left((O_{Q^*_j}\otimes\I^B)\,\omega_{k,m}\right)-t^*_jE_j(\psi)\Big) \leq \sum_j \epsilon_j.
\end{equation}
Thus, if $s_{k,m} > \sum_j \epsilon_j$, $\ket{\psi}$ cannot be converted into $\rho$ by LOCC. If the condition does not hold, the test does not tell us anything, in which case, for $k\geq3$, we can extract the decomposition from $\Phi_{k,m}$ and repair it as in Sec.~\ref{sec:extractingTheDecomposition}, and if the criterion Eq.~\eqref{eq:LOCCcriterion} holds, the state can be converted into $\rho$ by LOCC.
In this procedure, the grid errors $\epsilon_j$ need to be computed in a postprocessing step from the polynomials returned by the solver.

\subsection{Detection at finite level}
We impose $P_j\leq E_j$ on all spectra. 
Consider the decomposition $\{p_i, \dyad{\varphi_i}\}$ of $\rho$ that minimizes
\begin{equation}
    s^*=\max_t \sum_j t_j \left( \left(\sum_i p_i E_j(\varphi_i)\right)  - E_j(\psi) \right)
\end{equation}
over probability vectors $t$.
If $\ket{\psi}\overset{\text{LOCC}}{\nrightarrow}\rho$, at least one Vidal monotone $E_j$ will be exceeded by the decomposition and setting $t_j=1$, we get that $s^*>0$. 
We know $\widetilde{s}_{k,m} \leq s^*$ since we are relaxing the set of feasible $\omega_{k,m}$.
By App.~\ref{app:ProofCongNonExtHier}, as $k,m\to\infty$, the few copy marginals of $\Phi_{k,m}$ approach mixtures $\sum q_i\dyad{\phi_i}^{\otimes n}$ with $\{q_i, \dyad{\phi_i}\}$ a decomposition of $\rho$, and since by App.~\ref{app:LPconvergence} each $E_j$ can be approximated to arbitrary precision from below, $\widetilde{s}_{k,m} \to s^*$ as $k,m \to \infty$ and every non-convertible pair is eventually detected.
\section{Further analysis}
In this appendix, we provide additional numerical results for the linear entropy (Fig.~\ref{fig:namedLinearEntropy}) and the Vidal monotones (Fig.~\ref{fig:Vidal}).
% We compare our method of extracting upper bounds from the decomposition with a HJW based gradient descent method.
We also show how the LP adapts to different states by choosing different lower-bounding polynomials, which are tangent to the reference function at different points of the spectrum (Fig.~\ref{fig:polynomials}).


% Shared colors and styles of the figures. Every figure reads its data from figures/data/*.csv, written by
% experiments/make_figures.py (comma separated, header row, nan for missing values: pgfplots skips those points).
\definecolor{cA}{HTML}{2a78d6}\definecolor{cB}{HTML}{eb6834}\definecolor{cC}{HTML}{1baf7a}%
\definecolor{cD}{HTML}{4a3aa7}\definecolor{cE}{HTML}{c03b7a}\definecolor{cF}{HTML}{8a6d1f}\definecolor{cGray}{HTML}{8a8984}%
\definecolor{lvA}{HTML}{c6dbf5}\definecolor{lvB}{HTML}{93bbec}\definecolor{lvC}{HTML}{5e98e0}%
\definecolor{lvD}{HTML}{2a78d6}\definecolor{lvE}{HTML}{174a8c}\definecolor{lvF}{HTML}{1baf7a}%
%
% \namedbars{csv file}{ylabel}{title}: lower bounds of the named states at every level, divided by the reference `ref`
% (printed above each state at height `label_y`), with the two upper bounds as markers. Columns: see named_csv in
% make_figures.py. The state positions (`pos`) leave a gap at x = 10 between the PPT and the NPT states.
\providecommand{\namedbar}[3]{% fill, column, offset
  \addplot[ybar, bar width=0.13, bar shift=0pt, draw=none, fill=#1, area legend]
    table[x expr=\thisrow{pos}+#3, y expr=\thisrow{#2}/\thisrow{ref}] {\namedfile};}
\providecommand{\namedbars}[3]{%
  \def\namedfile{#1}%
  \begin{tikzpicture}
    \begin{axis}[width=\textwidth, height=0.29\textwidth, ymin=0, ymax=1.38, xmin=0.3, xmax=14.7,
      ylabel={#2}, ytick={0,0.25,0.5,0.75,1}, xtick=data, xticklabels from table={#1}{state},
      table/col sep=comma, xticklabel style={rotate=35, anchor=north east, font=\scriptsize, inner sep=1pt},
      tick label style={font=\footnotesize}, label style={font=\small}, title={#3}, title style={font=\small,
      yshift=8pt}, clip=false, legend style={font=\scriptsize, at={(0.5,-0.42)}, anchor=north, legend columns=4,
      draw=none, /tikz/every even column/.append style={column sep=6pt}}, legend cell align=left,
      ytick style={draw=none}, xtick style={draw=none}, ymajorgrids, grid style={cGray!25}]
      \addplot[draw=none, forget plot] table[x=pos, y expr=0] {#1};% the ticks (xtick=data)
      \namedbar{lvA}{lb12}{-0.325}\addlegendentry{LB $(1,2)$}
      \namedbar{lvB}{lb13}{-0.195}\addlegendentry{LB $(1,3)$}
      \namedbar{lvC}{lb22}{-0.065}\addlegendentry{LB $(2,2)$}
      \namedbar{lvD}{lb23}{0.065}\addlegendentry{LB $(2,3)$}
      \namedbar{lvE}{lb33}{0.195}\addlegendentry{LB $(3,3)$}
      \namedbar{lvF}{lb24}{0.325}\addlegendentry{LB $(2,4)$}
      \addplot[only marks, mark=diamond*, mark size=2.2pt, cB, mark options={draw=white, line width=0.4pt}]
        table[x expr=\thisrow{pos}-0.14, y expr=\thisrow{ub_extract}/\thisrow{ref}] {#1};
      \addlegendentry{UB: extracted ($k=m=3$)}
      \addplot[only marks, mark=x, mark size=2.4pt, black, thick]
        table[x expr=\thisrow{pos}+0.14, y expr=\thisrow{ub_hjw}/\thisrow{ref}] {#1};
      \addlegendentry{UB: HJW}
      \addplot[only marks, mark=none, nodes near coords={$\pgfplotspointmeta$}, point meta=explicit symbolic,
        nodes near coords style={anchor=south, font=\tiny, inner sep=1pt}, forget plot]
        table[x=pos, y=label_y, meta=ref_label] {#1};
      \draw[cGray] (axis cs:0.3,1) -- (axis cs:14.7,1);
      \draw[cGray, densely dashed] (axis cs:10,0) -- (axis cs:10,1.38);
      \node[font=\scriptsize, anchor=south] at (axis cs:5,1.38) {PPT (bound entangled)};
      \node[font=\scriptsize, anchor=south] at (axis cs:12.5,1.38) {NPT};
    \end{axis}
  \end{tikzpicture}}
%
\begin{figure*}
  \centering
  \namedbars{figures/data/named_SL.csv}{$S_{L,\mathrm{LB}}\,/\,\mathrm{UB}$}{}
  \caption{Lower bounds on the linear entropy for different PPT and NPT states. Bars: lower bounds from the simplified hierarchy at the levels $(k,m)$ shown in the legend. Markers: upper bounds from the decomposition extracted from the SDP solution at $k=m=3$ (orange) and from HJW gradient descent over decompositions (black). All values are divided by the best upper bound UB, printed above each state: the exact convex roof for two qubits (marked $^*$), and the lowest of the two upper bounds otherwise.}
  \label{fig:namedLinearEntropy}
\end{figure*}

% Shared colors and styles of the figures. Every figure reads its data from figures/data/*.csv, written by
% experiments/make_figures.py (comma separated, header row, nan for missing values: pgfplots skips those points).
\definecolor{cA}{HTML}{2a78d6}\definecolor{cB}{HTML}{eb6834}\definecolor{cC}{HTML}{1baf7a}%
\definecolor{cD}{HTML}{4a3aa7}\definecolor{cE}{HTML}{c03b7a}\definecolor{cF}{HTML}{8a6d1f}\definecolor{cGray}{HTML}{8a8984}%
\definecolor{lvA}{HTML}{c6dbf5}\definecolor{lvB}{HTML}{93bbec}\definecolor{lvC}{HTML}{5e98e0}%
\definecolor{lvD}{HTML}{2a78d6}\definecolor{lvE}{HTML}{174a8c}\definecolor{lvF}{HTML}{1baf7a}%
%
% \namedbars{csv file}{ylabel}{title}: lower bounds of the named states at every level, divided by the reference `ref`
% (printed above each state at height `label_y`), with the two upper bounds as markers. Columns: see named_csv in
% make_figures.py. The state positions (`pos`) leave a gap at x = 10 between the PPT and the NPT states.
\providecommand{\namedbar}[3]{% fill, column, offset
  \addplot[ybar, bar width=0.13, bar shift=0pt, draw=none, fill=#1, area legend]
    table[x expr=\thisrow{pos}+#3, y expr=\thisrow{#2}/\thisrow{ref}] {\namedfile};}
\providecommand{\namedbars}[3]{%
  \def\namedfile{#1}%
  \begin{tikzpicture}
    \begin{axis}[width=\textwidth, height=0.29\textwidth, ymin=0, ymax=1.38, xmin=0.3, xmax=14.7,
      ylabel={#2}, ytick={0,0.25,0.5,0.75,1}, xtick=data, xticklabels from table={#1}{state},
      table/col sep=comma, xticklabel style={rotate=35, anchor=north east, font=\scriptsize, inner sep=1pt},
      tick label style={font=\footnotesize}, label style={font=\small}, title={#3}, title style={font=\small,
      yshift=8pt}, clip=false, legend style={font=\scriptsize, at={(0.5,-0.42)}, anchor=north, legend columns=4,
      draw=none, /tikz/every even column/.append style={column sep=6pt}}, legend cell align=left,
      ytick style={draw=none}, xtick style={draw=none}, ymajorgrids, grid style={cGray!25}]
      \addplot[draw=none, forget plot] table[x=pos, y expr=0] {#1};% the ticks (xtick=data)
      \namedbar{lvA}{lb12}{-0.325}\addlegendentry{LB $(1,2)$}
      \namedbar{lvB}{lb13}{-0.195}\addlegendentry{LB $(1,3)$}
      \namedbar{lvC}{lb22}{-0.065}\addlegendentry{LB $(2,2)$}
      \namedbar{lvD}{lb23}{0.065}\addlegendentry{LB $(2,3)$}
      \namedbar{lvE}{lb33}{0.195}\addlegendentry{LB $(3,3)$}
      \namedbar{lvF}{lb24}{0.325}\addlegendentry{LB $(2,4)$}
      \addplot[only marks, mark=diamond*, mark size=2.2pt, cB, mark options={draw=white, line width=0.4pt}]
        table[x expr=\thisrow{pos}-0.14, y expr=\thisrow{ub_extract}/\thisrow{ref}] {#1};
      \addlegendentry{UB: extracted ($k=m=3$)}
      \addplot[only marks, mark=x, mark size=2.4pt, black, thick]
        table[x expr=\thisrow{pos}+0.14, y expr=\thisrow{ub_hjw}/\thisrow{ref}] {#1};
      \addlegendentry{UB: HJW}
      \addplot[only marks, mark=none, nodes near coords={$\pgfplotspointmeta$}, point meta=explicit symbolic,
        nodes near coords style={anchor=south, font=\tiny, inner sep=1pt}, forget plot]
        table[x=pos, y=label_y, meta=ref_label] {#1};
      \draw[cGray] (axis cs:0.3,1) -- (axis cs:14.7,1);
      \draw[cGray, densely dashed] (axis cs:10,0) -- (axis cs:10,1.38);
      \node[font=\scriptsize, anchor=south] at (axis cs:5,1.38) {PPT (bound entangled)};
      \node[font=\scriptsize, anchor=south] at (axis cs:12.5,1.38) {NPT};
    \end{axis}
  \end{tikzpicture}}
%
\begin{figure*}
  \centering
  \namedbars{figures/data/named_E1.csv}{$E_{1,\mathrm{LB}}\,/\,$ref}{Vidal monotone $E_1=1-\alpha_1$}
  \caption{Lower bounds on the Vidal monotone $E_1$ for different PPT and NPT states, from the SDP with embedded LP. Bars: lower bounds from the simplified hierarchy at the levels $(k,m)$ shown in the legend. Markers: upper bounds from the decomposition extracted from the SDP solution at $k=m=3$ (orange) and from HJW gradient descent over decompositions (black). All values are divided by the reference printed above each state: the exact convex roof for two qubits (marked $^*$), and the lowest of the two upper bounds otherwise.}
  \label{fig:Vidal}
\end{figure*}

% % Shared colors and styles of the figures. Every figure reads its data from figures/data/*.csv, written by
% experiments/make_figures.py (comma separated, header row, nan for missing values: pgfplots skips those points).
\definecolor{cA}{HTML}{2a78d6}\definecolor{cB}{HTML}{eb6834}\definecolor{cC}{HTML}{1baf7a}%
\definecolor{cD}{HTML}{4a3aa7}\definecolor{cE}{HTML}{c03b7a}\definecolor{cF}{HTML}{8a6d1f}\definecolor{cGray}{HTML}{8a8984}%
\definecolor{lvA}{HTML}{c6dbf5}\definecolor{lvB}{HTML}{93bbec}\definecolor{lvC}{HTML}{5e98e0}%
\definecolor{lvD}{HTML}{2a78d6}\definecolor{lvE}{HTML}{174a8c}\definecolor{lvF}{HTML}{1baf7a}%
%
% \namedbars{csv file}{ylabel}{title}: lower bounds of the named states at every level, divided by the reference `ref`
% (printed above each state at height `label_y`), with the two upper bounds as markers. Columns: see named_csv in
% make_figures.py. The state positions (`pos`) leave a gap at x = 10 between the PPT and the NPT states.
\providecommand{\namedbar}[3]{% fill, column, offset
  \addplot[ybar, bar width=0.13, bar shift=0pt, draw=none, fill=#1, area legend]
    table[x expr=\thisrow{pos}+#3, y expr=\thisrow{#2}/\thisrow{ref}] {\namedfile};}
\providecommand{\namedbars}[3]{%
  \def\namedfile{#1}%
  \begin{tikzpicture}
    \begin{axis}[width=\textwidth, height=0.29\textwidth, ymin=0, ymax=1.38, xmin=0.3, xmax=14.7,
      ylabel={#2}, ytick={0,0.25,0.5,0.75,1}, xtick=data, xticklabels from table={#1}{state},
      table/col sep=comma, xticklabel style={rotate=35, anchor=north east, font=\scriptsize, inner sep=1pt},
      tick label style={font=\footnotesize}, label style={font=\small}, title={#3}, title style={font=\small,
      yshift=8pt}, clip=false, legend style={font=\scriptsize, at={(0.5,-0.42)}, anchor=north, legend columns=4,
      draw=none, /tikz/every even column/.append style={column sep=6pt}}, legend cell align=left,
      ytick style={draw=none}, xtick style={draw=none}, ymajorgrids, grid style={cGray!25}]
      \addplot[draw=none, forget plot] table[x=pos, y expr=0] {#1};% the ticks (xtick=data)
      \namedbar{lvA}{lb12}{-0.325}\addlegendentry{LB $(1,2)$}
      \namedbar{lvB}{lb13}{-0.195}\addlegendentry{LB $(1,3)$}
      \namedbar{lvC}{lb22}{-0.065}\addlegendentry{LB $(2,2)$}
      \namedbar{lvD}{lb23}{0.065}\addlegendentry{LB $(2,3)$}
      \namedbar{lvE}{lb33}{0.195}\addlegendentry{LB $(3,3)$}
      \namedbar{lvF}{lb24}{0.325}\addlegendentry{LB $(2,4)$}
      \addplot[only marks, mark=diamond*, mark size=2.2pt, cB, mark options={draw=white, line width=0.4pt}]
        table[x expr=\thisrow{pos}-0.14, y expr=\thisrow{ub_extract}/\thisrow{ref}] {#1};
      \addlegendentry{UB: extracted ($k=m=3$)}
      \addplot[only marks, mark=x, mark size=2.4pt, black, thick]
        table[x expr=\thisrow{pos}+0.14, y expr=\thisrow{ub_hjw}/\thisrow{ref}] {#1};
      \addlegendentry{UB: HJW}
      \addplot[only marks, mark=none, nodes near coords={$\pgfplotspointmeta$}, point meta=explicit symbolic,
        nodes near coords style={anchor=south, font=\tiny, inner sep=1pt}, forget plot]
        table[x=pos, y=label_y, meta=ref_label] {#1};
      \draw[cGray] (axis cs:0.3,1) -- (axis cs:14.7,1);
      \draw[cGray, densely dashed] (axis cs:10,0) -- (axis cs:10,1.38);
      \node[font=\scriptsize, anchor=south] at (axis cs:5,1.38) {PPT (bound entangled)};
      \node[font=\scriptsize, anchor=south] at (axis cs:12.5,1.38) {NPT};
    \end{axis}
  \end{tikzpicture}}
%
\begin{figure}
  \centering
  \begin{tikzpicture}
    \begin{axis}[width=\columnwidth, height=0.62\columnwidth, ymode=log, ymin=1e-10, ymax=1,
      ylabel={$(S_L^{\mathrm{UB}}-S_{L,\mathrm{LB}_{3,3}})/S_{L,\mathrm{LB}_{3,3}}$},
      xtick=data, xticklabels from table={figures/data/upper_bounds_SL.csv}{state}, xmin=0.3, xmax=14.7,
      xticklabel style={rotate=50, anchor=north east, font=\tiny, inner sep=1pt}, tick label style={font=\footnotesize},
      label style={font=\small}, legend style={font=\scriptsize, at={(0.5,1.02)}, anchor=south, draw=none,
      legend columns=2, /tikz/every even column/.append style={column sep=6pt}}, legend cell align=left,
      ymajorgrids, grid style={cGray!25}, xtick style={draw=none}, table/col sep=comma]
      \def\file{figures/data/upper_bounds_SL.csv}
      \addplot[draw=none, forget plot] table[x=pos, y expr=1e-5] {\file};% the ticks (xtick=data)
      \addplot[only marks, mark=diamond*, mark size=2.2pt, cB, mark options={draw=white, line width=0.4pt}]
        table[x expr=\thisrow{pos}-0.15, y=gap_extract] {\file};\addlegendentry{UB: extracted ($k=m=3$)}
      \addplot[only marks, mark=x, mark size=2.4pt, black, thick]
        table[x expr=\thisrow{pos}+0.15, y=gap_hjw] {\file};\addlegendentry{UB: HJW}
      \draw[cGray, densely dashed] (axis cs:10,1e-10) -- (axis cs:10,1);
    \end{axis}
  \end{tikzpicture}
  \caption{
  Relative gap between the upper bounds and the SDP lower bounds. We show that the gap between the upper and lower bounds closes quite fast and the upper bounds extracted from the SDP solution are often tighter than the ones from HJW based gradient descent even at low levels.}
  \label{fig:upperBounds}
\end{figure}

% Shared colors and styles of the figures. Every figure reads its data from figures/data/*.csv, written by
% experiments/make_figures.py (comma separated, header row, nan for missing values: pgfplots skips those points).
\definecolor{cA}{HTML}{2a78d6}\definecolor{cB}{HTML}{eb6834}\definecolor{cC}{HTML}{1baf7a}%
\definecolor{cD}{HTML}{4a3aa7}\definecolor{cE}{HTML}{c03b7a}\definecolor{cF}{HTML}{8a6d1f}\definecolor{cGray}{HTML}{8a8984}%
\definecolor{lvA}{HTML}{c6dbf5}\definecolor{lvB}{HTML}{93bbec}\definecolor{lvC}{HTML}{5e98e0}%
\definecolor{lvD}{HTML}{2a78d6}\definecolor{lvE}{HTML}{174a8c}\definecolor{lvF}{HTML}{1baf7a}%
%
% \namedbars{csv file}{ylabel}{title}: lower bounds of the named states at every level, divided by the reference `ref`
% (printed above each state at height `label_y`), with the two upper bounds as markers. Columns: see named_csv in
% make_figures.py. The state positions (`pos`) leave a gap at x = 10 between the PPT and the NPT states.
\providecommand{\namedbar}[3]{% fill, column, offset
  \addplot[ybar, bar width=0.13, bar shift=0pt, draw=none, fill=#1, area legend]
    table[x expr=\thisrow{pos}+#3, y expr=\thisrow{#2}/\thisrow{ref}] {\namedfile};}
\providecommand{\namedbars}[3]{%
  \def\namedfile{#1}%
  \begin{tikzpicture}
    \begin{axis}[width=\textwidth, height=0.29\textwidth, ymin=0, ymax=1.38, xmin=0.3, xmax=14.7,
      ylabel={#2}, ytick={0,0.25,0.5,0.75,1}, xtick=data, xticklabels from table={#1}{state},
      table/col sep=comma, xticklabel style={rotate=35, anchor=north east, font=\scriptsize, inner sep=1pt},
      tick label style={font=\footnotesize}, label style={font=\small}, title={#3}, title style={font=\small,
      yshift=8pt}, clip=false, legend style={font=\scriptsize, at={(0.5,-0.42)}, anchor=north, legend columns=4,
      draw=none, /tikz/every even column/.append style={column sep=6pt}}, legend cell align=left,
      ytick style={draw=none}, xtick style={draw=none}, ymajorgrids, grid style={cGray!25}]
      \addplot[draw=none, forget plot] table[x=pos, y expr=0] {#1};% the ticks (xtick=data)
      \namedbar{lvA}{lb12}{-0.325}\addlegendentry{LB $(1,2)$}
      \namedbar{lvB}{lb13}{-0.195}\addlegendentry{LB $(1,3)$}
      \namedbar{lvC}{lb22}{-0.065}\addlegendentry{LB $(2,2)$}
      \namedbar{lvD}{lb23}{0.065}\addlegendentry{LB $(2,3)$}
      \namedbar{lvE}{lb33}{0.195}\addlegendentry{LB $(3,3)$}
      \namedbar{lvF}{lb24}{0.325}\addlegendentry{LB $(2,4)$}
      \addplot[only marks, mark=diamond*, mark size=2.2pt, cB, mark options={draw=white, line width=0.4pt}]
        table[x expr=\thisrow{pos}-0.14, y expr=\thisrow{ub_extract}/\thisrow{ref}] {#1};
      \addlegendentry{UB: extracted ($k=m=3$)}
      \addplot[only marks, mark=x, mark size=2.4pt, black, thick]
        table[x expr=\thisrow{pos}+0.14, y expr=\thisrow{ub_hjw}/\thisrow{ref}] {#1};
      \addlegendentry{UB: HJW}
      \addplot[only marks, mark=none, nodes near coords={$\pgfplotspointmeta$}, point meta=explicit symbolic,
        nodes near coords style={anchor=south, font=\tiny, inner sep=1pt}, forget plot]
        table[x=pos, y=label_y, meta=ref_label] {#1};
      \draw[cGray] (axis cs:0.3,1) -- (axis cs:14.7,1);
      \draw[cGray, densely dashed] (axis cs:10,0) -- (axis cs:10,1.38);
      \node[font=\scriptsize, anchor=south] at (axis cs:5,1.38) {PPT (bound entangled)};
      \node[font=\scriptsize, anchor=south] at (axis cs:12.5,1.38) {NPT};
    \end{axis}
  \end{tikzpicture}}
%
% \polypanel{measure column prefix (EF or E1)}: the optimal polynomial of each state, and the points where it touches F
% (the Schmidt coefficients of the Wootters decomposition).
\providecommand{\polypanel}[1]{%
  \addplot[cA, thick] table[x=alpha, y=werner2] {figures/data/polynomials_#1.csv};\addlegendentry{Werner $2\times2$}
  \addplot[cA, only marks, mark=*, mark size=1.8pt, forget plot]
    table[x=alpha_werner2, y=#1_werner2] {figures/data/polynomials_touching.csv};
  \addplot[cB, thick] table[x=alpha, y=gisin] {figures/data/polynomials_#1.csv};\addlegendentry{Gisin}
  \addplot[cB, only marks, mark=*, mark size=1.8pt, forget plot]
    table[x=alpha_gisin, y=#1_gisin] {figures/data/polynomials_touching.csv};
  \addplot[cC, thick] table[x=alpha, y=mems] {figures/data/polynomials_#1.csv};\addlegendentry{MEMS}
  \addplot[cC, only marks, mark=*, mark size=1.8pt, forget plot]
    table[x=alpha_mems, y=#1_mems] {figures/data/polynomials_touching.csv};
  \addplot[cD, thick] table[x=alpha, y=bellprod] {figures/data/polynomials_#1.csv};\addlegendentry{Bell + $\ket{01}$}
  \addplot[cD, only marks, mark=*, mark size=1.8pt, forget plot]
    table[x=alpha_bellprod, y=#1_bellprod] {figures/data/polynomials_touching.csv};}%
\begin{figure*}
  \centering
  \begin{tikzpicture}
    \begin{groupplot}[group style={group size=2 by 1, horizontal sep=1.6cm}, width=0.46\textwidth,
      height=0.32\textwidth, tick label style={font=\footnotesize}, label style={font=\small}, title style={font=\small},
      xlabel={$\alpha_1$}, ylabel={$F(\alpha)$, $P(\alpha)$}, xmin=0, xmax=1, table/col sep=comma,
      legend style={font=\scriptsize, at={(0.5,0.02)}, anchor=south, draw=none, fill=white, fill opacity=0.85,
      text opacity=1}, legend cell align=left]
    \nextgroupplot[title={(a) Entanglement of formation, level $(3,4)$}, ymin=-0.8, ymax=1.15]
      \addplot[cGray, dashed, forget plot] coordinates {(0,0) (1,0)};
      \addplot[black, very thick] table[x=alpha, y=F] {figures/data/polynomials_EF.csv};\addlegendentry{$S(\alpha)$}
      \polypanel{EF}
    \nextgroupplot[title={(b) Vidal monotone $E_1$, level $(3,3)$}, ymin=-0.4, ymax=0.62]
      \addplot[cGray, dashed, forget plot] coordinates {(0,0) (1,0)};
      \addplot[black, very thick] table[x=alpha, y=F] {figures/data/polynomials_E1.csv};\addlegendentry{$E_1(\alpha)$}
      \polypanel{E1}
    \end{groupplot}
  \end{tikzpicture}
  \caption{
  Polynomials $P\leq F$ extracted from the SDP with embedded LP for different $2$-qubit states. For each state, this is the polynomial that yields the best bound on the convex roof on the grid, at level $(3,4)$ for $E_F$ (a) and at level $(3,3)$ for $E_1$ (b). Here we show
  that the LP is adaptive depending on the state and there is not a single polynomial that yields the best bound on the convex roof for all states at the same level. The states used are $\{$ Werner $2\times2$: singlet weight $p=0.8$; Gisin: $\lambda=0.9$, $\theta=\pi/8$; MEMS: $C=0.8$; Bell + $\ket{01}$: $p=0.6$ $\}$. The polynomials are shown as a function of the entanglement spectrum $\a = (\a_1, \a_2) = (\a_1, 1-\a_1)$.}
  \label{fig:polynomials}
\end{figure*}


\section{Numerical implementation}\label{app:numerics}
All the SDPs are solved with the interior-point conic solver QICS~\cite{qics}, with its default tolerances.
The PPT entangled states from the figures come from the library \texttt{bound-entangled}~\cite{boundEntangled} and the NPT ones are built with \texttt{toqito}~\cite{toqito}. For the Horodecki $3\times3$ and $2\times4$ states we take the parameter $a=0.5$ in Figs.~\ref{fig:EoF}, \ref{fig:namedLinearEntropy} and~\ref{fig:Vidal}. The parameters of the NPT states are the ones given in the caption of Fig.~\ref{fig:polynomials}.

\subsection{Upper bounds by gradient descent}\label{app:HJW}
By the HJW theorem~\cite{HJW}, every decomposition of $\rho=\sum_i\lambda_i\dyad{e_i}$ 
% its laaaaate
into $M$ pure states is obtained
by mixing the vectors $\sqrt{\lambda_i}\ket{e_i}$ with an $M\times M$ unitary, so every unitary gives a decomposition of
$\rho$ and hence an upper bound on the convex roof. As in Ref.~\cite{Rothlisberger2009}, we minimize this upper bound by
gradient descent over the unitaries, with $M=r^2$ and several random starting points, and keep the best value.

% Shared colors and styles of the figures. Every figure reads its data from figures/data/*.csv, written by
% experiments/make_figures.py (comma separated, header row, nan for missing values: pgfplots skips those points).
\definecolor{cA}{HTML}{2a78d6}\definecolor{cB}{HTML}{eb6834}\definecolor{cC}{HTML}{1baf7a}%
\definecolor{cD}{HTML}{4a3aa7}\definecolor{cE}{HTML}{c03b7a}\definecolor{cF}{HTML}{8a6d1f}\definecolor{cGray}{HTML}{8a8984}%
\definecolor{lvA}{HTML}{c6dbf5}\definecolor{lvB}{HTML}{93bbec}\definecolor{lvC}{HTML}{5e98e0}%
\definecolor{lvD}{HTML}{2a78d6}\definecolor{lvE}{HTML}{174a8c}\definecolor{lvF}{HTML}{1baf7a}%
%
% \namedbars{csv file}{ylabel}{title}: lower bounds of the named states at every level, divided by the reference `ref`
% (printed above each state at height `label_y`), with the two upper bounds as markers. Columns: see named_csv in
% make_figures.py. The state positions (`pos`) leave a gap at x = 10 between the PPT and the NPT states.
\providecommand{\namedbar}[3]{% fill, column, offset
  \addplot[ybar, bar width=0.13, bar shift=0pt, draw=none, fill=#1, area legend]
    table[x expr=\thisrow{pos}+#3, y expr=\thisrow{#2}/\thisrow{ref}] {\namedfile};}
\providecommand{\namedbars}[3]{%
  \def\namedfile{#1}%
  \begin{tikzpicture}
    \begin{axis}[width=\textwidth, height=0.29\textwidth, ymin=0, ymax=1.38, xmin=0.3, xmax=14.7,
      ylabel={#2}, ytick={0,0.25,0.5,0.75,1}, xtick=data, xticklabels from table={#1}{state},
      table/col sep=comma, xticklabel style={rotate=35, anchor=north east, font=\scriptsize, inner sep=1pt},
      tick label style={font=\footnotesize}, label style={font=\small}, title={#3}, title style={font=\small,
      yshift=8pt}, clip=false, legend style={font=\scriptsize, at={(0.5,-0.42)}, anchor=north, legend columns=4,
      draw=none, /tikz/every even column/.append style={column sep=6pt}}, legend cell align=left,
      ytick style={draw=none}, xtick style={draw=none}, ymajorgrids, grid style={cGray!25}]
      \addplot[draw=none, forget plot] table[x=pos, y expr=0] {#1};% the ticks (xtick=data)
      \namedbar{lvA}{lb12}{-0.325}\addlegendentry{LB $(1,2)$}
      \namedbar{lvB}{lb13}{-0.195}\addlegendentry{LB $(1,3)$}
      \namedbar{lvC}{lb22}{-0.065}\addlegendentry{LB $(2,2)$}
      \namedbar{lvD}{lb23}{0.065}\addlegendentry{LB $(2,3)$}
      \namedbar{lvE}{lb33}{0.195}\addlegendentry{LB $(3,3)$}
      \namedbar{lvF}{lb24}{0.325}\addlegendentry{LB $(2,4)$}
      \addplot[only marks, mark=diamond*, mark size=2.2pt, cB, mark options={draw=white, line width=0.4pt}]
        table[x expr=\thisrow{pos}-0.14, y expr=\thisrow{ub_extract}/\thisrow{ref}] {#1};
      \addlegendentry{UB: extracted ($k=m=3$)}
      \addplot[only marks, mark=x, mark size=2.4pt, black, thick]
        table[x expr=\thisrow{pos}+0.14, y expr=\thisrow{ub_hjw}/\thisrow{ref}] {#1};
      \addlegendentry{UB: HJW}
      \addplot[only marks, mark=none, nodes near coords={$\pgfplotspointmeta$}, point meta=explicit symbolic,
        nodes near coords style={anchor=south, font=\tiny, inner sep=1pt}, forget plot]
        table[x=pos, y=label_y, meta=ref_label] {#1};
      \draw[cGray] (axis cs:0.3,1) -- (axis cs:14.7,1);
      \draw[cGray, densely dashed] (axis cs:10,0) -- (axis cs:10,1.38);
      \node[font=\scriptsize, anchor=south] at (axis cs:5,1.38) {PPT (bound entangled)};
      \node[font=\scriptsize, anchor=south] at (axis cs:12.5,1.38) {NPT};
    \end{axis}
  \end{tikzpicture}}
%
% \memcurve{column}{color}{mark}{legend}: one curve of memory.csv, broken between the groups of levels of equal degree m
% (pos 1-2: m = 2, pos 4-6: m = 3, pos 8-11: m = 4).
\providecommand{\memcurve}[4]{%
  \addplot[#2, thick, mark=#3, mark size=1.8pt, restrict x to domain=0:2.5] table[x=pos, y=#1] {\memfile};
  \addlegendentry{#4}
  \addplot[#2, thick, mark=#3, mark size=1.8pt, restrict x to domain=3.5:6.5, forget plot] table[x=pos, y=#1] {\memfile};
  \addplot[#2, thick, mark=#3, mark size=1.8pt, restrict x to domain=7.5:11.5, forget plot] table[x=pos, y=#1] {\memfile};}%
\providecommand{\fulllevel}[1]{\fill[cGray!15] (axis cs:#1-0.4,1e3) rectangle (axis cs:#1+0.4,1e16);}%
\begin{figure}
  \centering
  \def\memfile{figures/data/memory.csv}%
  \begin{tikzpicture}
    \begin{axis}[width=\columnwidth, height=0.8\columnwidth, ymode=log, ymin=1e3, ymax=1e16, xmin=0.4, xmax=11.6,
      xtick=data, xticklabels from table={figures/data/memory.csv}{level}, xlabel={level $(k,m)$},
      ytick={1e3,1e6,1e9,1e12,1e15}, yticklabels={1\,kB,1\,MB,1\,GB,1\,TB,1\,PB}, ylabel={memory},
      xticklabel style={rotate=45, anchor=north east, font=\scriptsize, inner sep=1pt},
      tick label style={font=\footnotesize}, label style={font=\small}, xtick style={draw=none},
      ymajorgrids, grid style={cGray!25}, clip=false, table/col sep=comma,
      legend style={font=\scriptsize, at={(0.5,-0.3)}, anchor=north, draw=none, legend columns=2,
      /tikz/every even column/.append style={column sep=4pt}}, legend cell align=left]
      \fulllevel{2}\fulllevel{6}\fulllevel{11}
      \addplot[draw=none, forget plot] table[x=pos, y expr=1e9] {\memfile};% the ticks (xtick=data)
      \draw[black, densely dashed] (axis cs:0.4,1.6e10) -- (axis cs:11.6,1.6e10)
        node[pos=0, anchor=south west, font=\scriptsize, inner sep=1pt] {16\,GB};
      \memcurve{mem_r4a2}{cA}{*}{$r=4$, $r_A=2$}
      \memcurve{mem_r4a3}{cB}{square*}{$r=4$, $r_A=3$}
      \memcurve{mem_r9a3}{cC}{triangle*}{$r=9$, $r_A=3$}
      \memcurve{mem_r4a4}{cD}{diamond*}{$r=4$, $r_A=4$}
      \memcurve{mem_r8a4}{cE}{pentagon*}{$r=8$, $r_A=4$}
      \memcurve{mem_r16a4}{cF}{star}{$r=16$, $r_A=4$}
      \node[font=\scriptsize, anchor=south] at (axis cs:1.5,1e16) {$m=2$};
      \node[font=\scriptsize, anchor=south] at (axis cs:5,1e16) {$m=3$};
      \node[font=\scriptsize, anchor=south] at (axis cs:9.5,1e16) {$m=4$};
    \end{axis}
  \end{tikzpicture}
  \caption{Memory needed by an interior-point solver for the bound at each level $(k,m)$, as a function of
  $r=\operatorname{rank}\rho$ and $r_A=\operatorname{rank}\rho_A$, for the full hierarchy [$k=m$, shaded,
  Eq.~\eqref{eq:SDPcroof1}] and the simplified one [$k<m$, Eq.~\eqref{eq:marginals_SDP}]. After the reductions of
  Appendix~\ref{app:SymRed}, the SDP only depends on $r$, $r_A$ and $(k,m)$; we count its number $n$ of free real
  variables (after eliminating the permutation-invariance equalities) and show the size $8n^2$ bytes of the dense
  Schur complement matrix, which dominates the memory of the solver for large problems. At equal degree $m$, the
  simplified levels with $k\geq2$ are cheaper than the full level only for states of high rank ($r=9$, $r_A=3$ and
  $r=16$, $r_A=4$); for low-rank states ($r=4$, $r_A=3$ and $r=8$, $r_A=4$) the full level $(m,m)$ needs less memory,
  since then $r<(k+1)r_A-k$ and a full copy costs less than a marginal copy. Dashed: the 16\,GB of a laptop.}
  \label{fig:memory}
\end{figure}


\end{widetext}

\end{document}
